\documentclass[preprint,10pt]{elsarticle}

\usepackage{amsmath,amssymb}
\usepackage{graphicx}
\usepackage{booktabs}
\usepackage{array}
\usepackage{xcolor}
\usepackage{algorithm}
\usepackage{algpseudocode}
\usepackage[section]{placeins}
\usepackage[colorlinks=true,linkcolor=blue,citecolor=blue,urlcolor=blue]{hyperref}

% ── figure palette (Okabe-Ito, color-blind safe; matches matplotlib exports) ──
\definecolor{iddink}{RGB}{30,58,84}
\definecolor{iddblue}{RGB}{0,114,178}
\definecolor{iddsky}{RGB}{86,180,233}
\definecolor{iddgreen}{RGB}{0,158,115}
\definecolor{iddamber}{RGB}{230,159,0}
\definecolor{iddamberdark}{RGB}{163,92,23}
\definecolor{iddgrey}{RGB}{143,143,143}

% ── AEI/Elsevier float discipline (best effort) ──────────────────────────────
% Floats are confined to their section (placeins [section]) with the class's
% standard separation (~12-15 pt). Best effort, not a guarantee: large tables
% and figures are placed with top preference and may drift a few pages from
% their first citation in the preprint build when deferred float pages flush.
\setlength{\floatsep}{12pt plus 2pt minus 2pt}
\setlength{\textfloatsep}{15pt plus 2pt minus 4pt}
\setlength{\intextsep}{12pt plus 2pt minus 2pt}
\setlength{\abovecaptionskip}{6pt}
\setlength{\belowcaptionskip}{2pt}
\renewcommand{\topfraction}{0.92}
\renewcommand{\dbltopfraction}{0.92}
\renewcommand{\textfraction}{0.01}
\renewcommand{\floatpagefraction}{0.85}
\makeatletter
\setlength{\@fptop}{0pt} % float pages: top-align content instead of vertical centring
\makeatother
\emergencystretch=1.5em

\journal{Advanced Engineering Informatics}

\begin{document}

\title{Auditable agentic diagnosis for industrial processes: evidence-graded competing-hypothesis reasoning evaluated against per-fault literature baselines}

\author[ustc]{Haozheng Zhang\corref{cor1}}
\cortext[cor1]{Corresponding author.}
\address[ustc]{Institute of Advanced Technology, University of Science and Technology of China, Hefei 230088, China}
\ead{zhanghaozheng@ustc.edu.cn}

\begin{frontmatter}
\begin{abstract}
Large language model (LLM) diagnosis of industrial processes promises training-free root-cause analysis, yet its conclusions are rarely auditable, and its evaluation scores sample-level classification where engineers ask scenario-level, mechanism-grounded questions. We present the Industrial Deep Diagnostic (IDD) pipeline: nine-stage agentic orchestration---deterministic statistics, competing-hypothesis reasoning, evidence grading (L1--L7), an anti-spurious-correlation filter, bounded-repair gates---with a machine-checked evaluation framework of provenance-labeled artifacts. On a 12-scenario benchmark spanning SKAB, Tennessee Eastman Process, and IndPenSim, IDD resolves 4/9 documented faults to a single root cause (Top-1 $44.4\%$---a provisional upper bound, since run-directory names can embed the documented mechanism; the joint worst-case floor is 0/9; identity-strict 3/9), and all nine name the true mechanism in their primary or surviving set (top-$k$ 9/9 at the keyword instrument's family level---8/9 under its stricter branch-localized reading). The five non-resolving cases return capped \mbox{\textsc{competing\_set}} verdicts naming the discriminating measurement; no control raises a false alarm (one of three blind). A same-model ablation bounds every accuracy claim: a bare call over the same blind digest keyword-scores 9/9 faults against the pipeline's 4/9 resolved (same-rule family-level comparison: 9/9 vs 9/9), so the pipeline's claimed margins are process qualities---auditable, capped-honest architecture---whose correctness anchoring by blinded expert adjudication remains the outstanding validation step. One seeded re-execution (a single repeat pair) reproduced verdict state, actuator, and mechanism class within two confidence points; the deterministic suite arms reproduce 45 recorded run outputs byte-identically (timestamps excluded). Throughout, \emph{trustworthy} denotes auditable and capped-honest, not verified correctness.
\end{abstract}

\begin{keyword}
Fault root-cause analysis \sep Reproducible benchmark \sep Large language models \sep Multi-agent systems \sep Competing hypotheses \sep Evidence grading \sep Process ontology
\end{keyword}

\end{frontmatter}

\section{Introduction}
\label{sec:intro}

Root-cause analysis (RCA) of industrial process anomalies is a knowledge-intensive activity: the engineer must decide \emph{which} of many correlated sensor deviations represents the true underlying cause, justify that judgment with evidence, and communicate it with a calibrated degree of confidence. Three long-standing obstacles keep that automation unsatisfactory in practice.

\textbf{First, single-LLM diagnosis is unstructured and unauditable.} Recent evaluations of LLM-based fault explanation on the Tennessee Eastman Process (TEP)~\cite{khan2024faultexplainer} show that large models can articulate plausible mechanisms, but their answers arrive as free-form text with no machine-checkable structure, no separation of facts from inference, and no mechanism that prevents a confident-sounding conclusion from resting on a spurious correlation. A recent review of agentic approaches to process monitoring and fault diagnosis concludes that evidence for LLM/multi-agent fault diagnosis remains scarce~\cite{agentic2026review}: the direction is promising but lacks engineering-grade pipelines.

\textbf{Second, classical data-driven fault detection and diagnosis (FDD) optimizes the wrong target.} Principal-component methods~\cite{kano2002comparison}, deep anomaly detectors~\cite{tuli2022tranad,pozdnyakov2024adversarial}, and their benchmarks~\cite{rieth2017additional} score \emph{per-sample} detection or classification~\cite{qin2012survey}. An engineer, however, asks a \emph{per-episode} question: given this anomalous run, what is the root cause, how confident are we, and what evidence supports it? A 99\%-accurate classifier that returns a fault label without mechanism, evidence, or calibration does not answer that question. The detection--diagnosis gap is not merely theoretical: classical detectability studies place the documented disturbances IDV3, IDV9 and IDV15 of the fifteen-fault TEP catalogue at the detection limit, with IDV4 weakly detected by the $T^2$ statistic~\cite{russell2000delay,chiang2001fault}, and the published LLM baseline's own PCA implementation leaves all four of these faults unscored~\cite{khan2024faultexplainer}.

\textbf{Third, correlation is not causation.} In recirculating process loops, almost every pair of variables is correlated through control action or material balance. An RCA method that treats strong correlation as causation will confidently report control-loop artifacts as root causes. Despite decades of causal-inference literature~\cite{yuan2014root,bauer2007finding}, contemporary LLM-diagnosis studies rarely include explicit anti-spurious safeguards.

This paper addresses the three obstacles with the Industrial Deep Diagnostic (IDD) pipeline and a benchmark built to evaluate it honestly. Our contributions are:

\begin{enumerate}
\item \textbf{A quality-gated agentic diagnosis pipeline} that couples deterministic statistical tooling with LLM reasoning under a competing-hypotheses protocol (Section~\ref{sec:method}), bounded repair loops, and automatic physical verification, so that every conclusion is decomposable into hypothesis, evidence, exclusion, and confidence.
\item \textbf{An evidence-grading and calibration scheme} (levels L1--L7 with gate propagation) that bounds reported confidence by the weakest supporting evidence level. In the corpus no \mbox{\textsc{determined}} verdict is wrong at the keyword instrument's family level (identity-strict: one, the batch-93 verdict; Section~\ref{sec:overall}), and in every scenario whose surviving mechanisms are not separable at the protocol's elimination thresholds the pipeline returns a capped \mbox{\textsc{competing\_set}} verdict that keeps the true mechanism ranked and names the discriminating measurement instead of forcing a conclusion (Section~\ref{sec:results}).
\item \textbf{A reproducible evaluation framework}---truth isolation, data fingerprinting, artifact-provenance labels, recorded execution traces, and a machine-checked artifact-consistency gate---in which a scenario is graded only if its complete artifact set was authored by gated pipeline stages and the execution proof passes; an internal integrity audit removed a scripted artifact-fabrication shortcut, and the benchmark was re-executed end to end under the tightened contract (Section~\ref{sec:framework}).
\item \textbf{A 12-scenario benchmark} on three public datasets (SKAB, TEP, IndPenSim) whose TEP subset carries scored per-fault baselines from FaultExplainer~\cite{khan2024faultexplainer} on four faults (the remaining two lie where its PCA screen excludes them). IDD resolves 4/9 faults to a single root cause---a provisional upper-bound reading under the disclosed identifier channel, pending the committed neutralized-identifier re-run (Section~\ref{sec:discussion})---while ranking the true mechanism in \emph{every} fault scenario's surviving set at the keyword instrument's family level (top-$k$ 9/9, a set-membership bound) and raising zero false alarms on the controls (one of three blind; the batch-93-excluded reading is 3/8); a seeded random-draw re-execution reproduced one scenario's verdict state and mechanism (Section~\ref{sec:stability}); and the graded runs executed as true per-stage sub-agent dispatches communicating only through workspace artifacts (Sections~\ref{sec:framework}, \ref{sec:results}--\ref{sec:comparison}).
\end{enumerate}

Throughout, \emph{trustworthy} is used in a delimited sense that the paper's evidence is expected to carry: an IDD conclusion is trustworthy to the extent that it is \emph{auditable}---every claim decomposes into hypotheses, evidence with levels, exclusions, and confidence, all machine-checkable from released artifacts---and \emph{capped-honest}: when the recording cannot separate the surviving mechanisms, the pipeline must say so and name the discriminating measurement rather than commit. Correctness is carried by the truth comparison of Section~\ref{sec:framework}, never by the label; what the pipeline adds is that its conclusions can be audited, reproduced, and escalated rather than merely read.

Section~\ref{sec:related} positions the work against data-driven FDD, LLM-based diagnosis, and multi-agent orchestration research. Sections~\ref{sec:arch} and \ref{sec:method} detail the architecture and methodology. Section~\ref{sec:framework} defines the evaluation framework. Sections~\ref{sec:setup}--\ref{sec:comparison} present the benchmark and comparisons, and Section~\ref{sec:discussion} discusses limitations and deployment prospects.

\section{Related Work}
\label{sec:related}

\subsection{Data-Driven Fault Detection and the Benchmark Gap}
Statistical process monitoring on the Tennessee Eastman Process established the dominant evaluation template~\cite{downs1993plant,qin2012survey}: principal-component analysis (PCA) and projection methods with Hotelling's $T^2$ (the model-subspace score) and the residual $Q$, or squared prediction error, statistic~\cite{kano2002comparison,yin2014review}, later extended by deep anomaly detectors~\cite{hartung2023deep,tuli2022tranad} and adversarial robustness studies~\cite{pozdnyakov2024adversarial}. Within the tripartite partition of fault detection and diagnosis into quantitative-model-based, qualitative-model-based, and process-history methods~\cite{venkatasubramanian2003review,venkatasubramanian2003part2,venkatasubramanian2003part3}, these families occupy the process-history branch, and the quantitative branch is anchored by model-based analytical redundancy and residual-based fault isolation~\cite{isermann2005model}. All of these branches score \emph{sample-level} detection or \emph{fault-class} accuracy. Critically, they do not evaluate root-cause \emph{mechanisms}: a classifier that labels a run ``IDV(4)'' provides the class identity but neither a causal mechanism, an evidence trail, nor a confidence statement. The classical detection ceiling (IDV3/9/15 at the limit, IDV4 weak) is documented~\cite{russell2000delay,chiang2001fault}. Benchmark methodology has drawn scrutiny: random segment splits inflate reported AUROC---by up to 15.3 percentage points of Macro-AUROC under segmentation-level leakage in Vieira et al.~\cite{vieira2026towards}, with a lineage formalized by Hendriks et al.~\cite{hendriks2022towards}---motivating our entity-level scenario splits (Section~\ref{sec:framework}). Scenario-level root-cause localization itself has a classical lineage that IDD must be read against: Granger-causal and transfer-entropy methods localize propagation directions in plant-wide oscillations~\cite{yuan2014root,bauer2007finding,yoon2001statistical,thornhill2007fault}, signed-directed-graph reasoning performs mechanism-chain elimination~\cite{iri1979algorithm}, reconstruction-based contribution analysis isolates root causes from propagated effects~\cite{dunia1996identification,alcala2009reconstruction}, and digital-twin practice (surveyed in~\cite{tao2019digital}) mirrors model-based diagnosis.

\subsection{LLM-Based Industrial Diagnosis}
FaultExplainer~\cite{khan2024faultexplainer} is the closest published system: it couples PCA-based detection with LLM-generated fault explanations on TEP. Its evaluation uses two prompt regimes. In the root-causes-included regime the prompt carries the documented root-cause list, the model outputs a top-3 candidate list, and alias matches are accepted (alias classes: IDV1/2/8, IDV3$\leftrightarrow$IDV9, IDV4$\leftrightarrow$IDV11$\leftrightarrow$IDV14, IDV5$\leftrightarrow$IDV12$\leftrightarrow$IDV15); here GPT-4o resolves 7/11 and o1-preview 9/11 of the 11 PCA-detectable faults (FaultExplainer's own PCA implementation excludes IDV3, IDV4, IDV9 and IDV15 from scoring). In its General-Reasoning regime---no candidate list, designed to simulate unseen faults---both models place the correct or a related cause in the top-3 for 8/11 faults under lenient ``related'' grading (FE's own text names faults 10 and 13 as failing yet reports 8/11; we transcribe the reported figure). Protocol differences therefore condition the comparison: (i) FE's outputs are top-3 lists with alias/related acceptance, whereas IDD commits to a single mechanism verdict scored by strict mechanism keywords; (ii) the two PCA-excluded faults in our subset (IDV3, IDV4) are scored by IDD; (iii) FE's alias policy grades a cooling-water-family answer as correct for IDV11 and IDV14 alike, whereas IDD separates the three forms of the same physical channel. Gong et al.~\cite{gong2026harnessing} orchestrate multiple LLM agents with difficulty-aware routing, but evaluate on a multiple-choice QA benchmark (FailureSensorIQ~\cite{constantinides2025failuresensoriq}) rather than on real time-series data.

Liu et al.~\cite{liu2025probing} combine scene graphs, fine-tuned Llama-3 and augmented reality for machine-tool maintenance, and causal-graph RCA has been studied for time-varying processes~\cite{liu2025spatial} and knowledge-map-based group decision settings~\cite{yue2023root}; a label-free evaluation scheme for diagnosis models has been proposed in the same venue~\cite{liu2024label}. Multimodal LLM frameworks extend the modality range: Alsaif et al.~\cite{alsaif2024multimodal} couple textual, tabular, and image modalities for Industry-4.0 fault detection and diagnosis over industrial fault records, and Merkelbach and Hasenauer~\cite{merkelbach2024multimodal} show that multimodal LLMs can draft physical models for automated diagnosis, within disclosed limits.
Recent work in this journal shows the same momentum and the same gaps.  CausalKGPT couples an industrial causal knowledge graph with an LLM for quality-problem cause analysis~\cite{zhou2024causalkgpt}; FD-LLM fine-tunes a large language model for equipment fault diagnosis~\cite{lin2025fdllm}; Wen et al.\ orchestrate LLMs for human--machine collaborative troubleshooting~\cite{wen2025troubleshooting}; and Xu et al.\ pair a deep digital twin with a vision-language model for multi-scenario fault diagnosis~\cite{xu2026digitaltwin}. Closest to our setting, the contemporaneous AgentRCA (July 2026) builds an agentic RCA system that maintains a hypothesis table over tool-grounded evidence, reporting 40.0\% Top-1 and 61.5\% Top-2 on its TEP evaluation set (it excludes IDV3/9/15 for the classical detectability reasons we discuss~\cite{wei2026agentrca}; its output is a ranked hypothesis list without calibrated confidence, and its evaluation, like FE's, stops short of an external truth-comparison harness, and its detectability exclusion also covers the unspecified-cause faults IDV16--20). Table~\ref{tab:related} summarizes the positioning: among published and contemporaneous systems, none combines real time-series inputs, scenario-level mechanism diagnosis with an evidence trail, cap-enforced three-state conclusions, and a machine-verifiable evaluation contract---the last two, in particular, are absent from every system cited above, with AgentRCA naming uncertainty communication as future work. The claim is scoped to industrial process diagnosis; the contemporaneous agentic-RCA line for software microservices targets a different artifact class (distributed traces and logs) and a different evidence regime~\cite{chen2023automatic}.

\begin{table}[!t]
\caption{Positioning Against the Closest Published Systems. Each row scores a published or contemporaneous system on real time-series input, scenario-level root-cause analysis, an evidence trail, cap-enforced confidence, and a reproducibility gate; IDD is the only system that combines all five. ``Repro.\ gate'' rates whether the system's evaluation is machine-replayed end to end (IDD: internal to our own harness, a boundary Table~\ref{tab:trustboundary} bounds); PCA/deep FDD baselines are ``partial'' because detection rates are reproducible but diagnosis claims are not.}
\label{tab:related}
\centering
\footnotesize
\resizebox{\textwidth}{!}{
\begin{tabular}{p{3.1cm}p{1.5cm}p{2.5cm}p{1.7cm}cc}
\toprule
System & Input & Scenario-level RCA & Evidence trail & Cap-enforced & Repro. gate \\
\midrule
FaultExplainer~\cite{khan2024faultexplainer} & TEP series & partial (LLM text) & partial & no & no \\
AgentRCA~\cite{wei2026agentrca} & TEP series & yes (ranked list) & partial (traces) & no & no \\
Gong et al.~\cite{gong2026harnessing} & MCQA only & no (QA acc.) & no & no & no \\
PCA / deep FDD (e.g., TranAD~\cite{tuli2022tranad}) & series & no (labels) & no & n/a & partial \\
Causal RCA (e.g.,~\cite{yue2023root}) & series & partial (graph) & partial & no & no \\
\textbf{IDD} & \textbf{series} & \textbf{yes (3-state)} & \textbf{yes (L1--L7)} & \textbf{yes} & \textbf{yes (internal)} \\
\bottomrule
\end{tabular}
}
\end{table}

\subsection{Multi-Agent Orchestration Frameworks}
Generic agent frameworks (AutoGen~\cite{wu2023autogen} and commercial derivatives; ReAct-style reasoning--acting loops~\cite{yao2023react}) provide message-passing and role prompting but leave audit, evidence policy, and evaluation contracts to the application. IDD's artifact-only communication is a blackboard architecture with schema-checked writes; recent failure taxonomies for multi-agent LLM systems \cite{cemri2025mast} catalogue the failure classes its gates address case by case (fabricated artifacts $\to$ C4, inter-agent misalignment $\to$ CP-5/6). Our position, following the practice of the reviewed agentic-FDD literature~\cite{agentic2026review}, is that industrial diagnosis requires these as \emph{first-class, machine-checked} pipeline elements rather than prompt conventions; Section~\ref{sec:arch} describes how IDD enforces them through checkpoints, bounded repair loops, and an event-log execution proof.

\subsection{Trust, Automation Levels, and Intellectual Ancestry}\label{sec:trust}
IDD's design decisions sit on well-mapped ground, and we credit the map. The three-state verdict is a level-of-automation and trust-calibration decision---what the system concludes versus what it escalates to humans---in the sense of Parasuraman et al.'s types and levels of automation~\cite{parasuraman2000model}---a line descending from Sheridan's early function-allocation analysis~\cite{sheridan1978human}---and Lee \& See's appropriate-reliance account of trust~\cite{lee2004trust}: a capped \mbox{\textsc{competing\_set}} implements calibrated transfer in their sense, and the audit trail's effect on its human reader is two-sided, sustaining appropriate reliance where it is checked while inviting automation bias where it is not~\cite{parasuraman2010complacency,bainbridge1983ironies}. The competing-hypotheses discipline descends from Heuer's Analysis of Competing Hypotheses~\cite{heuer1999psychology}, the four-condition anti-spurious filter from Hill's causation criteria~\cite{hill1965environment}, and the abstention option from Chow's recognition-with-reject tradeoff~\cite{chow1970optimum}, whose statistical counterpart is modern selective prediction~\cite{angelopoulos2023conformal}; the L1--L7 grounding check doubles as a faithfulness proxy in the sense of rigorous XAI evaluation~\cite{doshivelez2017rigorous}. The named-discriminating-measurement exit is the passive counterpart of \emph{active} fault diagnosis~\cite{scott2014input}: IDD inherits the passive setting of plant recording and converts its residual ambiguity into an instrumentation request. The elimination calculus itself descends from model-based diagnosis---Reiter's consistency-based theory~\cite{reiter1987theory} and its multiple-fault extension~\cite{dekleer1987diagnosing}; IDD supplies the consistency oracle with per-recording statistics and, where consistency cannot be reached, returns the three-state verdict instead of forcing a single fault. The ontology-asset layer continues the OntoCAPE lineage~\cite{morbach2009ontocape} and, behind it, the ISO~15926 plant-data modeling tradition~\cite{iso15926}. IDD's contribution with respect to this ancestry is a domain-specific, machine-checked instantiation of these frameworks inside an agentic pipeline---not their invention.

\section{System Architecture}
\label{sec:arch}

\subsection{Execution Architecture}
The following terms recur throughout. A \emph{skill package} is a versioned folder holding an agent's protocol checklist, JSON Schemas, and validation scripts; a \emph{specialist agent} is an independently spawned session that executes one skill package; an \emph{engine} is the CLI/RPC harness that runs such a session. Fig.~\ref{fig:arch} shows the skill execution architecture as five planes, of which the audit plane is machine-checked. Client surfaces (the web application, CLI drivers, and a CI harness) submit jobs through the engine layer: fourteen interchangeable CLI/RPC engines behind one uniform event contract, with pre-flight probing and typed error codes. The orchestrating main agent holds the domain knowledge as \emph{skill packages} and dispatches each pipeline stage by passing absolute paths (\texttt{SKILL\_PATH}, \texttt{RUN\_DIR}); sub-agents communicate only through workspace artifacts, never through conversation context.

Read against the released software, the layer map needs three qualifications. First, the package contains 18 skill packages executed by 14 specialized agents, but the \emph{diagnostic path} reported here (Steps 2--8.5) dispatches eight specialist sub-agents---context-builder, data-processor, diagnostician, judge, report-reviewer, reporter, html-visualizer, html-reviewer---with the VLM visual analyzer dispatched by the data-processor inside its Phase~5.5; the remaining packages form optional layers outside the benchmarked path (Section~\ref{sec:arch-layers}). Second, a skill package is not a prompt alone: every dispatched package carries the protocol checklist (\texttt{references/agent-protocol.md}) and ships the JSON Schemas and validation scripts that machine-check its artifacts before the stage may complete---the structural difference from prompting a model directly. Third, the engine layer gives the pipeline harness portability: all fourteen engines implement one start/stream/close contract, so an engine substitution is a configuration change rather than a re-implementation---pre-registered as an ablation axis, not yet exercised. Every stage transition appends a structured event to the audit plane, which the log checker and the finalization gate reduce to a single machine verdict; grading (Section~\ref{sec:framework}) reads only artifacts that survived this plane.

\subsection{Knowledge, Web, and Enhancement Layers}\label{sec:arch-layers}
Beyond the diagnostic path sit three further layers. They complete the description of the released system, but none altered a benchmark result: how each was exercised, degraded, or left idle is recorded in the run events.

\emph{Semantic-knowledge layer.} A retrieval-augmented generation microservice~\cite{lewis2020rag} (FastAPI + ChromaDB, \texttt{rag-retrieval-engine/}) supplies semantic understanding for ontology construction and is the only sanctioned source of level-L6 external evidence; the orchestrator probes it with a three-second health check at Step~0. When it is unavailable, the ontology builder falls back to the shipped \texttt{parameter\_to\_physics.json} mechanism map plus the data's own self-description, and the degradation is recorded in the run events. In the benchmark corpus the channel was never available: all thirteen runs record a failed probe and build every ontology through the documented fallback (recorded per run as a dedicated probe-failure file or a \texttt{rag\_validated:false} ontology flag); the TEP ontologies ground their variable semantics in the published process description, cross-checked against the data. No L6-tagged item exists in the released corpus, so no verdict rests on external evidence. Because the fallback ontology is the knowledge artifact every downstream stage actually consumed, it is worth showing what it contains (SKAB valve scenario, translated from the released bilingual artifact under \texttt{01\_ontology} in the run directory; abridged): the equipment register grounds each component in plant physics---centrifugal pump, ``operating point = intersection of pump and loop-resistance curves''; drive motor, ``current follows pump shaft load torque'' (derived $V\!\cdot\!I\approx226$\,W); throttle valve, ``regulates the loop resistance coefficient; position unmeasured---listed as the key discriminating uncertainty''---and a confounder register records the load-side and voltage confounders the diagnosis must exclude. The artifact is schema-validated and machine-checkable, and its variable semantics are what the ontology gate (CP-2/3) verifies before the statistical stage may run.

\emph{Ontology asset store.} Ontologies are published to a portable asset store (\texttt{ontology\_store.mjs publish}) keyed by a schema fingerprint; a later scenario whose variable schema matches is resolved by deterministic fast-reuse, which copies the validated ontology verbatim into the run directory in seconds instead of re-deriving it. All thirteen benchmark ontologies were published to the store on completion; none was \emph{consumed} by fast-reuse in this corpus, because the benchmark configuration pins the ontology step to full rebuilds (the reuse path is exercised by the production harness outside this benchmark). The store is deterministic and therefore does not affect the benchmark's truth-isolation contract (it reuses variable semantics, never scenario verdicts).

\emph{Integration surface.} The web application couples an Express backend with a Vue~3/Vite frontend and an SSE/WebSocket live stream. Its diagnosis endpoint accepts an engine selection and returns typed pre-flight errors for unknown or unusable engines, with probing and the real spawn sharing one code path. The optional enhancement layer (E0--E8, deterministic scripts orchestrated by a dedicated agent) post-processes completed runs for reporting depth and lies outside the benchmarked path.

\begin{figure}[!t]
\centering
\includegraphics[width=\textwidth,height=0.8\textheight,keepaspectratio]{figures/fig_architecture.png}
\caption{Skill execution architecture (diagram authored in draw.io; the \texttt{.drawio} source and the export script are released with the paper). Fourteen interchangeable engines expose one uniform event contract to the orchestrating agent, which dispatches the diagnostic path's eight specialist sub-agents from 18 skill packages; sub-agents communicate only through workspace artifacts, and the audit plane's log checker and finalization gate produce the execution proof the benchmark grader consumes.}
\label{fig:arch}
\end{figure}

\subsection{Pipeline Overview}
IDD organizes diagnosis as a nine-stage pipeline executed by the specialized agents backed by the skill packages (Fig.~\ref{fig:pipeline}; the benchmark's execution form---true per-stage sub-agent dispatch---is disclosed in Section~\ref{sec:framework}). The main orchestrator performs setup and data ingestion (Step~0/1) with adaptive preprocessing; a context-builder agent constructs a process ontology with variable semantics, physical principles, and confounder annotations (Step~2, checkpoint CP-2/3); a data-processor agent runs the deterministic statistical stage (per-variable anomaly profiling, cross-domain correlation, lag-compensated cross-correlation, batch integrity checks), producing a machine-readable analysis conclusion (Step~3, CP-4); a vision--language-model (VLM) agent reads rendered charts as visual evidence (pipeline Step~3.5 executes inside the data-processor agent as its internal Phase~5.5; in this deployment it executed its documented non-visual fallback in every run, Section~\ref{sec:setup}). The diagnostician agent then performs competing-hypothesis reasoning over the statistical and ontological evidence (Step~4, CP-5). Two parallel auditors, a ten-criterion judge and a physical pre-auditor, gate the conclusion (Step~5a/5b); the reporter assembles a pyramid-structured report (Step~6); a final physical audit must return ENDORSED before the HTML visualization and review (Steps~8/8.5, CP-9 at the review gate) and the evidence-closure finalization (Step~9).

\begin{table}[!t]
\caption{Module contracts of the diagnostic path: every stage, the agent that authors its artifacts, the artifacts it consumes and produces (filenames as released; all paths are relative to the run directory), and its checkpoint. The table lists each stage's contracted deliverables; runs additionally hold auxiliary stage-authored files. Checkpoint predicates: Table~\ref{tab:checkpoints}.}
\label{tab:modules}
\centering
\scriptsize\setlength{\tabcolsep}{3.2pt}\renewcommand{\arraystretch}{0.86}
\resizebox{\textwidth}{!}{%
\begin{tabular}{@{}>{\raggedright\arraybackslash}p{1.8cm}>{\raggedright\arraybackslash}p{1.7cm}>{\raggedright\arraybackslash}p{3.0cm}>{\raggedright\arraybackslash}p{4.9cm}c@{}}
\toprule
Stage & Agent & Consumes & Produces & Gate \\
\midrule
Setup/inspect (0/1) & main agent & scenario CSV, user context & \texttt{00\_input/input\_manifest.json}, \texttt{user\_context.json}, \texttt{run\_config.json} & CP-1 \\
Ontology (2) & context-builder & \texttt{input\_manifest.json}, RAG digest, process description & \texttt{01\_ontology/ontology.json} (+\texttt{00\_input/\allowbreak clarification\_needed.json}) & CP-2/3 \\
Statistics (3) & data-processor & \texttt{ontology.json}, cleaned data & \texttt{02\_processed/\allowbreak data\_analysis\_conclusion.json} (mandatory handoff), \texttt{validate\_report.json}, \texttt{03\_figures/} & CP-4 \\
Visual evidence (3.5) & vlm-visual-analyzer (inside data-processor Phase~5.5) & \texttt{03\_figures/}, \texttt{image\_captions.json} & \texttt{03\_figures/visual\_analysis.json} & --- \\
Diagnosis (4) & diagnostician & analysis conclusion, ontology, digests & \texttt{04\_diagnostics/\{diagnosis, evidence, confidence, reasoning\_chain\}.json} & CP-5 \\
Judge $\|$ pre-audit (5a/5b) & judge $\|$~report-reviewer & \texttt{04\_diagnostics/*} & \texttt{05\_review/judge\_feedback.json} $\|$ \texttt{05\_review/optimizer\_preflight.md} & CP-6 \\
Report (6) & reporter & all upstream artifacts & \texttt{report.md}, \texttt{run\_summary.json} & CP-7 \\
Final audit (7) & report-reviewer & \texttt{report.md}, artifacts & \texttt{optimizer.md} containing \mbox{ENDORSED} & CP-8 \\
HTML build (8) & html-visualizer & artifacts, design system & \texttt{render\_manifest.json} $\rightarrow$ \texttt{diagnostic-report.html} $\rightarrow$ \texttt{html\_selfcheck.json} & --- \\
HTML review (8.5) & html-reviewer & page + manifest & \texttt{05\_review/html\_review.json} (\texttt{verdict=pass}) & CP-9 \\
Finalize (9) & deterministic script & all of the above & \texttt{pipeline\_finalize\_report.json} (PASS), \texttt{evidence\_closure\_report.json} & --- \\
\bottomrule
\end{tabular}%
}
\end{table}

\begin{table}[!t]
\caption{Checkpoint gates CP-1\dots CP-9. Every predicate is mechanical---file existence, JSON-schema validity, or a numeric lower bound---so the gates can be re-run by a reader without domain judgment; a failing predicate blocks the transition and routes to the stated repair action.}
\label{tab:checkpoints}
\centering
\scriptsize\setlength{\tabcolsep}{4pt}
\resizebox{\textwidth}{!}{%
\begin{tabular}{@{}lcp{7.2cm}p{3.4cm}@{}}
\toprule
CP & Transition & Machine predicate & Failure action \\
\midrule
CP-1 & 1$\to$2 & \texttt{input\_manifest.json} + \texttt{user\_context.json} + \texttt{run\_config.json} exist & back to Step~0 \\
CP-2 & 2$\to$2.5 & \texttt{ontology.json} $\ge$1\,KB and schema-valid & re-run ontology builder \\
CP-3 & 2.5$\to$3 & \texttt{clarification\_status} $\in$ \{AUTO\_RESOLVED, USER\_CONFIRMED\} & resolve clarification \\
CP-4 & 3$\to$4 & \texttt{data\_analysis\_conclusion.json} present and generated plots $>0$ & re-run data-processor \\
CP-5 & 4$\to$5 & four diagnosis outputs schema-valid and quality-checked & repair diagnosis \\
CP-6 & 5$\to$6 & judge repair summary written; pre-audit reports no \mbox{FATAL} finding & repair (best-of-3) \\
CP-7 & 6$\to$7 & \texttt{report.md} + \texttt{run\_summary.json} present & re-run reporter \\
CP-8 & 7$\to$8 & \texttt{optimizer.md} contains ENDORSED & repair loop \\
CP-9 & 8.5 & \texttt{diagnostic-report.html} $\ge$5\,120\,B and review verdict \texttt{pass} & re-run html-visualizer \\
\bottomrule
\end{tabular}%
}
\end{table}

\textbf{Why nine stages.} Three rules govern the decomposition; the model-controlled ablation (Section~\ref{sec:ablation}) tests the first of them directly---the deterministic/semantic split---while the data-order and gated-audit rules rest on the design argument that follows. Deterministic work is pushed down to scripts: statistics, hashing, schema checks, and the finalization inventory run as versioned code so they are byte-reproducible, and only semantic judgment (ontology, hypotheses, causal framing, prose) is left to agents---a split that also keeps the deterministic layer defensible once a bare model call is known to be accurate (Section~\ref{sec:ablation}). Data description precedes diagnosis: the ontology and the statistics digest are produced before any hypothesis is generated, which is what forces the diagnose-only-what-the-data-supports behavior observed in the capped verdicts. Audit and delivery are gated last and separately: two independent auditors (judge and physical auditor, the pipeline's only mandatory parallel step) review the reasoning before a report exists, and the HTML deliverable is built manifest-first and reviewed independently, so the final artifacts cannot diverge from the audited verdict. The stage boundaries themselves are the artifact handoffs in Table~\ref{tab:modules}, and the mechanical predicates of the checkpoints are listed in Table~\ref{tab:checkpoints}.

\begin{figure}[!t]
\centering
\includegraphics[width=\textwidth,height=0.7\textheight,keepaspectratio]{figures/fig_pipeline.png}
\caption{Diagnostic pipeline flow for one scenario (diagram authored in draw.io; the \texttt{.drawio} source and the export script are released with the paper). Solid arrows: strictly sequential stages, each gated by a machine-validated checkpoint (blue tags); the judge and the physical pre-audit (Steps 5a$\,\|\,$5b) are the only mandatory parallel stages, drawn side by side feeding one gate. The dashed amber edges are the bounded repair loops: a sub-threshold judge score re-enters the gated diagnosis stages under the bounded-repair counters of Section~\ref{sec:method}, and a final audit without endorsement returns the report to Step~6---never an unbounded cycle. The second row reads right-to-left (serpentine flow); the return arrow at the row boundary marks the direction. Green boxes mark the hard gates: the final audit must return ENDORSED (CP-8), the HTML review gate must pass (CP-9), and the finalization gate must return \textsc{pass} (Step 9) before the deliverables exist. The VLM visual stage (Step 3.5, executed inside Step 3) ran its documented non-visual fallback in every reported run (Section~\ref{sec:setup}).}
\label{fig:pipeline}
\end{figure}

\subsection{Execution Layer: Interchangeable Engines With a Uniform Contract}
Reasoning agents are executed by any of fourteen interchangeable engines---the Claude Agent SDK, several RPC and ACP dialects, and a set of one-shot CLIs---behind a single start/stream/close contract realized as a capability table (\texttt{engines.mjs}) plus first-party wiring entries in the service dispatch table. Engine availability is probed and validated before execution with typed error codes (\texttt{400} unknown engine, \texttt{409 HARNESS\_UNAVAILABLE} unusable engine); a scripted in-process mock engine exercises the full pipeline in CI without external dependencies. Because every engine emits the same standardized event stream, the benchmark harness consumes any compliant engine's stream identically; diagnostic quality across engines is an empirical ablation question, and engine substitution remains a designed but not-yet-exercised ablation axis (Section~\ref{sec:discussion}).

\subsection{Audit Infrastructure}
Every stage transition appends a structured event to \texttt{.pipeline\_\allowbreak events.\allowbreak jsonl}; a log checker validates the start and completion records of every agent against the required contract, and a finalization gate aggregates inventory, schema, closure, judge-cross-audit, and event-archive checks into a single PASS/FAIL verdict. A run that fails any critical check cannot be graded, which is the property the benchmark relies on: \emph{no result exists without an execution proof}. The gate is exercised, not decorative, and it has produced one substantive diagnosis-layer intervention: the released records show the round-1 judge failing TEP IDV3 at 75 with two blocking issues, repaired to a round-2 pass at 90 under the declared scope---the audit plane's single demonstrated catch of a substantive diagnosis defect, which we report as the exception it is rather than as a rate. The remaining interventions were report-layer: the final physical audit caught two report-layer defects (a mis-attributed row reference and a residual over-claim in a run summary), repaired and re-audited; the grading rubric (Section~\ref{sec:setup}) records calibration and grounding deductions in three further runs; and several runs additionally bounced stage artifacts for schema or path-normalization issues before passing. Every rejection was repaired by the artifact's owning stage and re-gated to PASS; no verdict text changed in any of these report-layer repairs. Events are pipeline-internal records, not externally signed attestations; the resulting trust boundary and its limits are stated in Section~\ref{sec:discussion}.

\section{Diagnostic Methodology}
\label{sec:method}

\subsection{Problem Formulation}
A scenario supplies a recorded window $D_S$ of $p$ process variables and, optionally, user context; the task is to produce a verdict $v \in \{\textsc{det},\textsc{cs},\textsc{nd}\}$ with a mechanism-level root cause, a confidence $c \in [0,1]$, and an evidence chain. The diagnostician forms a hypothesis set $H = \{h_1,\dots,h_k\}$, $k \ge 3$, where each $h_i$ carries a mechanism class $m_i$, a set $E_i$ of evidence items, and a contradiction set $X_i$. Two definitions govern the protocol. The \emph{evidence level} of a hypothesis is the weakest level in its chain, $\ell(h_i) = \min_{e \in E_i} \mathrm{level}(e)$ with $\mathrm{level}(e) \in \{L1,\dots,L7\}$ (Section~\ref{sec:method-grading})---a conclusion cannot be stronger than its weakest citation. The \emph{confidence ceiling} is the minimum of the five-factor assessment and every protocol cap applicable to the verdict state,
\begin{equation}
c \le \bar{c}(v) = \min\bigl(F,\ \mathcal{C}(v)\bigr),
\label{eq:ceiling}
\end{equation}
where $F$ is the five-factor assessment and the cap function $\mathcal{C}(v)$ evaluates to the smallest applicable protocol cap: $0.70$ for \textsc{cs}, tightened to $0.65$ for its indistinguishable subclass, and $0.50$ whenever the anti-oscillation rule (code name \texttt{PARAM\_AMBIGUITY}) is active---all three are documented in the agent-side protocol. The agent-side protocol documents no \textsc{det} cap, but the machine grading gate flags---rather than blocks---confidence above $0.90$ for every non-\textsc{cs} state and above $0.70$ for \textsc{cs} (calibration check R4 of the released rubric; the two 0.91/0.93 verdicts of Section~\ref{sec:overall} are exactly this flag firing), so the effective cap set $\mathcal{C}$ includes $0.90$ for \textsc{det} at the gate (the gate's flat non-CS $0.90$ ceiling also covers \textsc{nd}). Reason-coded per-run ceilings may still bind; the released schema enumerates the reasons: indistinguishable competing set, unsorted-data lag claims, product-specific-only evidence, all-criteria-verified, direct measurement. For example, the canonical IDV7 run binds its \mbox{\textsc{determined}} verdict at $0.78$ under a reason-coded ceiling of $0.80$ (recorded reason \texttt{ALL\_CRITERIA\_VERIFIED}, noting that the root variable, header supply pressure, is unobservable, so the verdict is capped at the justified-inference level). Elimination is quantitative: $h_i$ is eliminated when its contradictions outweigh its support, $|X_i| > |E_i|$ with each item weighted by its evidence level (weights in the released protocol; the machine-checked floor is $\ge$2 eliminations, exclusion confidence $\ge 0.90$, and a citation per elimination), and every elimination records a revival condition. The three-state verdict is then the decision function
\begin{equation}
v = \begin{cases}
\textsc{det}, & |H_{\mathrm{surv}}| = 1 \ \text{and every elimination is evidenced},\\
\textsc{cs}, & |H_{\mathrm{surv}}| \ge 2 \ \text{and no channel separates the survivors},\\
\textsc{nd}, & \ell(h^\star) < L5 \ \text{or a required measurement is absent},
\end{cases}
\label{eq:verdict}
\end{equation}
with the reported confidence satisfying $c \le \bar{c}(v)$ (Eq.~\eqref{eq:ceiling}). The observed 0.55--0.63 band of capped verdicts reflects this INDISTINGUISHABLE subclass. The five-factor assessment is the protocol's confidence decomposition over statistical strength, physical plausibility, temporal evidence, absence of confounds, and symptom completeness (weighted 25/25/20/20/10, totalling 100, in the released \texttt{confidence.json}); the three states are protocol terms throughout this paper and carry no everyday meaning: \textsc{det} asserts a single evidenced mechanism, \textsc{cs} asserts that the present data cannot separate the survivors and names the measurement that would, and \textsc{nd} declares the evidence insufficient. Algorithm~\ref{alg:pipeline} states the whole procedure as executed, with deterministic (script) and semantic (agent) steps marked; the ablation of Section~\ref{sec:ablation} is the empirical test of this split.

\begin{algorithm}[!t]
\caption{IDD diagnosis of scenario $S$ (script = deterministic code, agent = semantic reasoning; gates as in Table~\ref{tab:checkpoints})}
\label{alg:pipeline}
\begin{algorithmic}[1]
\State \textbf{script:} fingerprint $D_S$, compute statistics digest, write blind brief \Comment{Step 0/1}
\State \textbf{agent:} $O \leftarrow$ ontology $(\text{brief})$ or \textbf{script}-side fast-reuse from the asset store \Comment{Step 2, CP-2/3}
\State \textbf{agent:} $A \leftarrow$ data-processor $(D_S, O)$; plots, validation report \Comment{Step 3, CP-4}
\State \textbf{agent:} form $H$, attach $(E_i, X_i)$ with evidence levels \Comment{Step 4}
\State \textbf{agent:} eliminate while $|X_i| > |E_i|$; compute $v$ by Eq.~\eqref{eq:verdict}
\State $c \leftarrow \min(\text{five-factor assessment}, \bar{c}(v))$ \Comment{Eq.~\eqref{eq:ceiling}}
\If{judge score $< 90$} \textbf{snapshot} best round; repair in declared scope ($\le 3$ rounds, global cap 5)
\EndIf \Comment{CP-6}
\State \textbf{agent:} report; \textbf{verify:} physical audit must return ENDORSED \Comment{Step 6/7, CP-7/8}
\State \textbf{agent:} HTML manifest-first, then independent review \Comment{Step 8/8.5, CP-9}
\State \textbf{script:} finalize---the grader's 14-artifact completeness set, schemas, closure, event log $\Rightarrow$ PASS; else the run cannot be graded (the Step-9 finalizer separately executes its own $\approx$45-check inventory)
\end{algorithmic}
\end{algorithm}

\subsection{Competing-Hypotheses Protocol}
The diagnostician agent must generate at least three candidate hypotheses per fault scenario (the protocol's intellectual ancestry is credited in Section~\ref{sec:trust}), assign each a mechanism class (e.g., fluid restriction, cavitation, rotor dynamics, feed-composition step, cooling disturbance, sensor drift), collect for each an evidence list with explicit levels, identify contradictions, and eliminate hypotheses whose contradictions are stronger than their support. The protocol admits exactly three conclusion states: \mbox{\textsc{determined}} (a single surviving hypothesis, bounded by the five-factor assessment; reason-coded per-run ceilings may further bind it), \mbox{\textsc{competing\_set}} ($\ge$2 survivors that the data cannot discriminate, with confidence capped at 0.70; observed caps in this corpus are 0.55--0.63 and the discriminating measurement is named), and \mbox{\textsc{needs\_data}} (evidence insufficient; missing data itemised). A fourth, implicit state---normal operation with no confirmed fault---applies to control scenarios. The elimination contract is quantitative: every elimination must cite the contradicting evidence, and every surviving chain must cite evidence numbers present in the statistical digests.

\subsection{Evidence Grading and Gate Propagation}\label{sec:method-grading}
Every piece of evidence carries a level L1--L7 (L1 direct measurement, L2 operating documents, L3 statistical analysis, L4 visual evidence, L5 domain ontology/mechanism, L6 external reference, L7 unsupported assumption). The ordering is a policy choice: direct measurement outranks statistical processing because a sensor reading cannot inherit a processing error, while visual evidence is graded below statistics because chart reading is itself an inference. The scenario confidence is bounded by the weakest level contributing to the surviving chain, and the quality gate flags, for deduction, conclusions whose confidence exceeds that bound. This converts ``hallucination control'' from a prompt instruction into a checkable arithmetic constraint.

\subsection{Anti-Spurious-Correlation Filter}
Four conditions must hold before a correlation may be cited as causal support: (i)~temporal precedence of cause over effect, verified by lag-compensated cross-correlation; (ii)~statistical significance surviving detrending, Simpson-stratification checks, and multiple-testing correction; (iii)~a physical mechanism annotated in the ontology connecting the variables; and (iv)~absence of contradictions under leave-one-out outlier sensitivity analysis, with production-state (steady/batch-phase) filtering applied before correlation. Correlations failing the filter are recorded but explicitly barred from hypothesis support; definitional pairs (e.g., valve position and the level it controls, $r\!=\!1.0$) are classified as control-loop artifacts. In the benchmark this filter is the reason control-loop correlations are not reported as root causes, and the reason normal batches with impulsive recipe-dosing spikes ($|z|$ up to 16.2) are graded as healthy rather than faulty.

\subsection{Automatic Physical Verification}
After the conclusion is drafted, a physics checker recomputes quantitative feasibility (mass/energy balances, thermal and vibration magnitude envelopes) from the cleaned data, ontology, and anomaly report, producing a machine-readable \texttt{physical\_checks} verdict. When no automated check applies to a scenario's column semantics, the pipeline falls back to an explicit L1--L5 manual-verification memo; silent omission is not permitted.

\subsection{Quality Gates, Repair Loops, and Calibration}
A ten-criterion judge---an LLM-as-judge, heir to that instrument's known biases (self-preference, verbosity \cite{zheng2023judging,panickssery2024selfpref})---scores the diagnosis (threshold 90); a parallel physical pre-audit reviews mechanism chains; failures trigger bounded repair (per-fault rerun $\le$3, global cap 5, with anti-oscillation confidence capping at 0.50). The final physical audit must return ENDORSED. If the review-repair budget exhausts without an endorsement, the run proceeds best-effort and the finalization gate records the missing endorsement as a critical failure, so the run is a pipeline failure rather than a gradeable result. Ceiling compliance is then auditable end-to-end: every scenario grading records whether the reported confidence respects the protocol ceiling, and \emph{overconfidence} (\mbox{\textsc{determined}} \emph{and} wrong) is a first-class error metric.

The released governance keeps a stochastic reasoner bounded by design. Each repair round snapshots the best-scoring attempt (\texttt{best\_round} copy of the diagnosis set); the second and third rounds run in a declared \emph{repair scope} that recomputes only the implicated artifacts and marks the rest as carried over, so a fix cannot silently rewrite evidence that was not found defective (auditable from the declared repair scopes and carried-over markers in the released logs). The loop never halts the run on score alone: if the threshold is still unmet after the budget, the pipeline delivers a best-effort report with a gate-enforced capped verdict (and the cap recorded), because failing to deliver is worse than delivering with a disclosed confidence. The protocol adds two machine-checkable disciplines. The \emph{data-truth mandate} requires every number in a diagnosis to be reproducible from the raw data (chain items carry derived markers and a NOT-APPLICABLE fallback), and the red-light list makes the principal failure modes auditable: no forced single verdict from a competing set, no global correlation cited as causal evidence, no conclusion without an evidence-level tag, no HTML outside the manifest. Violations are what the judge is instructed to flag, which is how the audit plane of Section~\ref{sec:arch} remains a check rather than a claim.

\section{A Reproducible Evaluation Framework}
\label{sec:framework}

Section~\ref{sec:method}'s gates are only as credible as the evaluation that scores their outputs, so the evaluation contract is fixed before any result is produced: IDD's benchmark framework (released as the project's \texttt{experience/} evidence repository and \texttt{docs/benchmark/} protocol) is designed around five contracts, each machine-enforced:

\textbf{(C1) Truth isolation.} Scenario definitions carry ground truth and grading keywords, but the pipeline's visible inputs contain only the process description and the data. Ground truth never enters the inference context; the grader owns it exclusively. The agent-visible brief itself carries exactly three fields---user question, process description, and a truth-withheld flag; dataset, role, expected-verdict-set, and grading metadata remain on the grader's side of the released bundle. A leakage sentinel machine-checks the contract from both directions: no ground-truth string may appear in any agent-visible brief, and no fault-vocabulary may appear in the run directories' user-context files. Scenario definitions---including the grading keywords and the allowed verdict set of every scenario---are versioned in the released repository with their edit history. The keyword sets used to grade the current corpus are the calibrated set pinned at commit \texttt{221d061}---over-generic tokens removed, mechanism synonyms added---then frozen before any current-era execution. The calibration is truth-adjacent in two ways that Section~\ref{sec:overall} quantifies: the ``v1-era corpus'' it was tuned against denotes the pre-discipline execution batch of the same twelve scenarios (run directories dated before the 2026-09-14 anti-oscillation and confidence-cap revision, retained with the released edit history), and the published per-fault baselines contributed the mechanism synonyms. The sentinel machine-checks briefs and user-context files, not run-directory names; the identifier channel is the residual disclosure channel and is treated in Section~\ref{sec:discussion}. These keys are author-designed (Section~\ref{sec:discussion}): a hit certifies mechanism-family agreement with the documented truth, not identity, and one scenario's set includes a generic token.

\textbf{(C2) Data fingerprinting.} Every input file is hashed (SHA-256) into a dataset manifest; the reproducibility gate re-verifies all 153 input fingerprints and the leakage sentinel returns a PASS verdict before any metric is accepted.

\textbf{(C3) Stage-authored artifacts with provenance.} Each pipeline stage writes its own machine-checkable artifacts: the ontology, the statistical analysis conclusion, the four-piece diagnosis, the judge's ten-criterion feedback, the auditors' ENDORSED verdicts, the report, the page model (\texttt{render\_manifest.json}) behind the HTML deliverable, and the reviewer's independent verdict. Every scored verdict carries an explicit provenance label (\texttt{judge-agent-10-criteria}, \texttt{auditor-agent}, \texttt{html-reviewer-agent}), so downstream consumers can distinguish independently authored judgments from rendered ones.

\textbf{(C4) Gate-graded scoring.} The grader accepts a scenario only if its run directory contains the complete artifact set of an executed pipeline (the fourteen items in Table~\ref{tab:artifacts}), the event log passes the log checker, and the finalization gate returns PASS. Otherwise the scenario is reported as \emph{not executed}---a pipeline failure, not a benchmark result. The gate is exercised in practice: an earlier protocol revision allowed a scripted shortcut that expanded a single reasoning note into a full artifact set (complete with synthesised agent events); the internal integrity audit flagged the fabrication, the shortcut was deleted from the codebase, and the entire benchmark was re-executed through the real pipeline. The grader's completeness gate makes the shortcut permanently unreplayable.

\textbf{(C5) Artifact-consistency gate.} A verifier recomputes all aggregate metrics from the per-scenario gradings, re-checks dataset hashes, verifies case coverage, and requires an execution proof per run; any drift fails the evaluation. The gate verifies the internal consistency of the recorded execution---recomputed aggregates, hashes, coverage, proofs---and is deliberately not an independent re-execution of the reasoning layer; the released report carries the resulting internally-consistent verdict.

\textbf{Execution form.} The canonical benchmark runs reported here were executed in the production dispatch form: an orchestrating agent dispatched each pipeline stage to an \emph{independent sub-agent session} through the harness's sub-agent spawn tool, passing absolute paths (\texttt{SKILL\_PATH}, \texttt{RUN\_DIR}); sub-agents communicated only through workspace artifacts, never through conversation context, and each authored its own artifacts under its own skill protocol. The pipeline's one mandatory parallel step (judge and pre-auditor, Steps 5a/5b) executed as two concurrent sessions. The orchestrator's role was limited to dispatch, gate verification, and event bookkeeping---it authored no pipeline artifact and no agent event. The benchmark driver itself retains the anti-fabrication division of labour: its scripts verify and score; where artifacts were missing it printed an execution contract and refused to grade. The event logs record every agent start/completion with declared outputs, and the log checker validates ordering and declared files per run. This form is the property that makes the audit plane meaningful: had stages shared one session, the artifact contracts and gates would be unchanged, but the independence between the authoring stage and its reviewers would be weaker than the production design intends. The independence is session- and protocol-level, not model-level: every stage runs on the one model family disclosed in Section~\ref{sec:setup}, so correlated failures between a generator and its reviewer remain possible.

Table~\ref{tab:artifacts} lists the graded artifact set the completeness gate enforces; the same fourteen filenames appear in the released run directories, so ``the complete artifact set'' is a checkable list rather than a phrase. One boundary should be read explicitly: the contracts' machine predicates are existence, schema, and threshold checks, whereas the epistemic constraints---evidence-level honesty, elimination substance, physical plausibility---are enforced by the reviewer stages under their protocols and audited arithmetically (ceiling compliance, grounding), not by those predicates; Table~\ref{tab:trustboundary} makes the resulting split between mechanism and discipline explicit.

\begin{table}[!t]
\caption{The graded artifact set enforced by the completeness gate (C4). A run missing any item is \emph{not executed} by the grader---a pipeline failure, not a benchmark result.}
\label{tab:artifacts}
\centering
\scriptsize\setlength{\tabcolsep}{4pt}
\resizebox{\textwidth}{!}{%
\begin{tabular}{@{}clp{6.6cm}@{}}
\toprule
\# & Run path & Role \\
\midrule
1 & \texttt{01\_ontology/ontology.json} & process ontology (variables, semantics, mechanisms) \\
2 & \texttt{02\_processed/data\_analysis\_conclusion.json} & statistical handoff to the diagnostician \\
3 & \texttt{04\_diagnostics/diagnosis.json} & verdict, hypotheses, eliminations \\
4 & \texttt{04\_diagnostics/evidence.json} & evidence inventory with numeric evidence ranks (1--7, the L1--L7 levels) \\
5 & \texttt{04\_diagnostics/confidence.json} & per-hypothesis confidence and ceiling audit \\
6 & \texttt{04\_diagnostics/reasoning\_chain.json} & R1--R8 reasoning chain with falsification conditions \\
7 & \texttt{05\_review/judge\_feedback.json} & judge sub-agent ten-criterion verdict \\
8 & \texttt{05\_review/optimizer\_preflight.md} & pre-report physical audit \\
9 & \texttt{05\_review/html\_review.json} & independent HTML review (verdict=pass) \\
10 & \texttt{report.md} & pyramid report \\
11 & \texttt{run\_summary.json} & machine-readable run summary (consistency-checked) \\
12 & \texttt{render\_manifest.json} & page model, written before the HTML \\
13 & \texttt{html\_selfcheck.json} & HTML self-check against the manifest \\
14 & \texttt{diagnostic-report.html} & HTML deliverable (size re-check rejects files below 5\,120\,B) \\
\bottomrule
\end{tabular}%
}
\end{table}

\textbf{Reproducibility semantics.} The framework distinguishes two layers explicitly. The deterministic statistical stage is byte-reproducible: identical inputs and code regenerate identical digests, and the released dataset fingerprints pin the inputs. The reasoning layer is reproducible at the \emph{procedure} level: re-executing the protocol re-launches the same gated pipeline on the same blind inputs and re-derives conclusions under the same gates and truth comparison; wording may vary between runs, but every verdict remains bound to the same evidence discipline. Reproduction claims in this paper therefore distinguish the two layers, and bit-level reproducibility is claimed only for the deterministic stage.

Scenario splits follow the entity-level discipline recommended by recent benchmark studies~\cite{vieira2026towards,hendriks2022towards}: each experimental batch/file appears in exactly one scenario, controls are paired with faults per dataset, and no segment-level splitting is used anywhere.

\section{Experimental Setup}
\label{sec:setup}

\subsection{Datasets and Scenarios}
We assemble 12 scenarios from three public benchmarks, listed in Table~\ref{tab:bench}. \textbf{SKAB}~\cite{skab2020} is a closed water-circulation testbed with vibration, current, pressure, temperature, and flow channels at 1\,Hz. Its two fault scenarios are inlet-valve throttling (valve1: the official experiment description names closing of the valve at the pump's flow inlet, i.e.\ suction-line restriction~\cite{skab2020}) and two-phase-flow cavitation drawn from the ``other'' group (other/13: the official description names two-phase flow supply to the pump inlet; the mechanism word is withheld from the scenario description); the anomaly-free control completes the set. \textbf{TEP}~\cite{downs1993plant,lyman1995plant,rieth2017additional} (IDV$n$ denotes entry $n$ of the documented disturbance catalogue; the simulator's revision history is documented by Bathelt et al.~\cite{bathelt2015revision}) contributes six documented fault scenarios: IDV1 feed-composition step, IDV3 D-feed-temperature step (classically hard to detect), IDV4 reactor-cooling-water step, IDV7 C-header pressure loss, IDV11 random cooling-water variation, and IDV14 cooling-valve stiction, plus the fault-free baseline as control. The TEP scenario briefs declare the anomaly-onset window (sample 161, idx 160 zero-based), so detection is out of scope for these scenarios and the task is RCA within a localized window; the IndPenSim briefs declare no onset. \textbf{IndPenSim}~\cite{goldrick2015development,goldrick2019modern} is a 100{,}000\,L penicillin fermentation simulator; its dataset documentation lists batches 91--100 as fault batches that ``contain faults resulting in process deviations''~\cite{goldrick2019modern}, and the selected scenario is batch 93. The dataset's own truth channel is deliberately mechanism-free: its fault-reference flag marks 101 of the 1,050 prepared rows as faulted (flag $=1$ from 70\,h into the batch), encoding fault/no-fault only---the simulation script's fault menu (aeration, back pressure, substrate feed, \dots) is a per-run parameter not published per batch---so the actuator-side deviation description (out-of-band pH and temperature excursions driven by the acid-dosing and heating-water flows) is the benchmark case definition's, not the dataset's and not this paper's diagnosis---and a case definition that may itself under-describe the batch disturbance---with one normal batch as control; the pipeline's cooling-loop signature for batch 93 is evaluated against this documented deviation in Section~\ref{sec:results}. 
Ground truth derives from verifiable sources only: SKAB's official label file, the established TEP cause table~\cite{chiang2001fault}, and IndPenSim's faulted-batch markers. Scenario selection followed a documentation rule rather than random sampling: every scenario corresponds to a publicly documented fault episode, chosen so that each dataset contributes mechanisms from distinct families plus a paired control; the consequences of non-random selection are discussed in Section~\ref{sec:discussion}. 
The control briefs mirror the fault briefs' structure but differ in what they disclose: all twelve briefs share the same anomaly-presupposing user question (``diagnose the root cause of this anomaly''); what differs is the process description: the SKAB control's is identical to its fault sibling's except that the fault brief announces an anomaly (onset withheld), while the TEP and IndPenSim control briefs state that the segment is the fault-free baseline---just as the corresponding fault briefs state that an anomaly exists. Only the SKAB control is therefore blind: the zero-false-alarm outcome gives one blind data point ($n{=}1$) that the pipeline does not fabricate a fault even when the question presupposes one and only the process description withholds the anomaly, and the SKAB control additionally exercises blind discrimination (its strong raw correlations are refuted, Section~\ref{sec:calibration}).

Three design decisions shaped the subset. First, the TEP subset is chosen so that \emph{every fault maps onto a published per-fault outcome} from FaultExplainer~\cite{khan2024faultexplainer}, enabling scenario-granular comparison rather than aggregate-only comparison: four faults (IDV1/7/11/14) carry scored baselines, and IDD additionally scores the two faults that screen excludes, IDV3 and IDV4. Second, IDV3 sits at the classical PCA detection limit and IDV4 is weakly detected~\cite{russell2000delay,chiang2001fault}---both lie where classical detection and the published baseline stop---so they probe root-cause reasoning at and beyond that limit. Third, IDV4, IDV11 and IDV14 perturb the same physical channel (reactor cooling water) with three different dynamic forms (step, random variation, valve stiction); separating them requires reading the statistics of each recording, not recalling a fault label, which partially mitigates prompt-knowledge contamination.

\begin{table}[!t]
\caption{Benchmark Composition. Twelve scenarios from three public datasets: each dataset contributes its documented fault scenarios plus one anomaly-free control, with ground truth drawn from official label files, the established TEP cause table, and IndPenSim's faulted-batch markers.}
\label{tab:bench}
\centering
\footnotesize
\resizebox{\textwidth}{!}{
\begin{tabular}{lccp{7.2cm}}
\toprule
Dataset & Fault & Control & Truth source \\
\midrule
SKAB & 2 & 1 & official \texttt{ground\_truth.json}; mechanism word withheld from the scenario brief for the ``other''-group case \\
TEP & 6 & 1 & Downs--Vogel; Chiang et al.~\cite{downs1993plant,chiang2001fault} \\
IndPenSim & 1 & 1 & faulted-batch markers (batch 93) \\
\midrule
Total & 9 & 3 &---\\
\bottomrule
\end{tabular}
}
\end{table}

\subsection{Metrics}
For fault scenarios we report \textbf{resolved Top-1} (a \mbox{\textsc{determined}} verdict naming the true mechanism; coincides with the FDD community's correct-diagnosis-rate, CDR), \textbf{top-$k$} (the true mechanism is named within the primary finding or the surviving-hypothesis set that the grader reads; because the protocol mandates at least three generated hypotheses, this approximates candidate-list membership, and for a single-survivor \mbox{\textsc{determined}} verdict it coincides with resolved Top-1; reported alongside, not instead of, resolved Top-1), \textbf{ceiling compliance} (reported confidence within the protocol ceiling of Section~\ref{sec:method}), \textbf{verdict-type compliance} (conclusion type within the scenario's allowed set; a \mbox{\textsc{determined}}-and-wrong verdict counts as \emph{overconfidence}), and a \textbf{protocol-compliance rubric score}---a deterministic, machine-computed checklist (100 points) computed from the pipeline artifacts themselves and measuring conformance to this paper's own protocol rather than correctness (see Section~\ref{sec:discussion}). Top-$k$ is structurally favored when several scenarios in a family share keyword tokens (the three reactor-cooling TEP faults share ``reactor cooling''); read it as family-level membership, not specific-form identification. The strictest, branch-localized reading would count IDV3 as a miss (its D/E-feed-temperature survivor sits on the stripper-feed branch while the documented truth is the reactor-feed D stream), giving top-$k$ 8/9 under that reading---both readings are reported. For control scenarios we report \textbf{pass} (no confirmed fault) and \textbf{false alarms}. Proportions carry Wilson 95\% confidence intervals at $z{=}1.96$ (the released scorer's constant). We do not report sample-level F1/AUC: they answer a different question and are used only as reference context (comparison protocol below).

The rubric's seven checks are computed strictly from the stage-authored run artifacts: \textbf{R1} artifact integrity (ontology, analysis conclusion, and the schema-valid diagnosis four-piece; 15 points), \textbf{R2} hypothesis structure ($\ge$3 hypotheses, $\ge$2 eliminations each citing substantive contradictory evidence; 15), \textbf{R3} evidence grounding (every surviving chain cites at least one number present in the scenario's own evidence brief; 20; fault scenarios---the check has no fault-relevant targets in a control brief, and the one control deduction under it is recorded as a check mis-specification, not a run defect; future corpora skip R3 for control scenarios by rule rather than by post-hoc exception), \textbf{R4} calibration (confidence within the protocol ceiling; conclusion type within the scenario's allowed set; 15), \textbf{R5} falsifiability (a surviving hypothesis names the observation that would revise it; 10), \textbf{R6} physics verification (automated check valid or explicit manual memo; 10), \textbf{R7} execution proof (event log + finalization PASS + HTML review gate; 15). The pipeline's internal ten-criterion judge remains the in-pipeline gate (Section~\ref{sec:method})---its scores are recorded with provenance and reported as execution-quality evidence---but it is \emph{not} a benchmark metric: benchmark quality scoring is fully mechanical. The rubric's scope is limited in two ways. First, it is a conformance instrument with limited discriminative power: nine of the twelve scenarios score at ceiling because the structural checks cannot fail for any run that has passed the gates---correctness claims rest on the truth comparison, never on the rubric. Second, ``strict mechanism-keyword hit'' is not mechanism equivalence: the keyword sets are disclosed per scenario, one of them (the IndPenSim process-deviation scenario) includes the generic token ``deviation'', and a hit certifies mechanism-family agreement, not identity. The consequences are quantified in Section~\ref{sec:results}, including an identity-strict Top-1 count that discounts the batch-93 verdict.\footnote{Artifact-field semantics for auditors: the released \texttt{metrics.json} folds verdict-type compliance into its \texttt{calibrated} flag, and \texttt{rubric.json} records the SKAB-valve type miss under its calibration-check key (the failing payload is the allowed-set violation described in Section~\ref{sec:calibration}); Table~\ref{tab:overall} reports calibration and verdict-type compliance as separate metrics.}

\subsection{Comparison Protocol}
We distinguish two comparison registers. \textbf{Register A (same-task direct comparison)}: methods that output one root-cause conclusion per scenario on the same dataset family; the published FaultExplainer per-fault numbers fall here, at scenario granularity. \textbf{Register B (reference only)}: sample-level detection/classification results; they are reported for context and must not be ranked against our scenario-level metrics. Register A comparisons inherit two protocol differences: FaultExplainer supplies candidate root causes in its prompt and accepts alias matches, whereas IDD uses no candidates and scores mechanism keywords strictly; and FaultExplainer scores only the 11 PCA-detectable faults, whereas IDD additionally scores the two faults that screen excludes, IDV3/IDV4. Because the scenario sets differ, no significance test is claimed across systems.

\subsection{Implementation}
The deterministic stages run as versioned scripts (streaming CSV inspection, anomaly profiling, correlation with lag-compensated cross-correlation functions (CCF), validation with Simpson/detrending/outlier-sensitivity/multiple-testing checks, physics checking via the shared scientific Python environment). The reasoning stages execute the full nine-stage agent pipeline under their skill protocols (Section~\ref{sec:arch}); every stage transition is recorded in the event log. The twelve canonical scenarios were executed as concurrent batches under one orchestrating agent (per-run event spans reflect orchestration batching, not isolated agent time); the clean per-scenario figure is the fully sequential random-draw re-execution (Section~\ref{sec:stability}; $\approx$2.0\,h end to end). The deterministic stage takes seconds per scenario. The entire chain (prepare, briefs, pipeline-execution verification, gate-graded scoring, rubric, aggregation, artifact-consistency gate, report) is driven by a single entry point with fail-fast stage gates.\medskip

\textbf{Benchmark configuration disclosure.} All thirteen runs reported here (twelve canonical scenarios plus one random-draw re-execution) were executed on one harness: the ZCode agent CLI (the CLI-driver class of the engine layer, dispatching each stage to an independent sub-agent session; execution form in Section~\ref{sec:framework}) with the GLM model family~\cite{glm2024chatglm} deployed under the provider's coding-plan configuration (snapshot September 2026). The harness does not expose temperature or top-p control, so provider defaults applied. The vision stage had no vision-capable model in this deployment and executed its documented non-visual degradation in every run (script-signed skeleton visual analysis with caption-based metadata fallback). The execution window was 17--18 September 2026, with one graded execution per scenario plus the drawn re-execution of TEP IDV11. Repair accounting: twelve of the thirteen runs closed with zero diagnosis-layer repair rounds; the exception (TEP IDV3, the hardest fault) spent one judge-driven targeted repair round, and two further scenarios required report-layer corrections after the final audit (Section~\ref{sec:arch})---every repair re-gated to PASS with no verdict text changed. Base-model capability is a confound for any cross-system comparison: the protocol-versus-baseline contrast of Section~\ref{sec:comparison} holds the prompt regime constant but not the model. The pre-registered ablation matrix targeting this attribution question has been partially executed (Section~\ref{sec:ablation}): a single LLM prompted with the identical blinded statistical digest, a classical PCA baseline on all twelve scenarios, and a FaultExplainer-protocol replication on the same model family. The engine-substitution arm and the filter-disabled arm remain future work. The 18 skill packages (the prompts executed by each agent) are released with the repository; \ref{app:artifacts} indexes the released per-scenario artifacts.

\section{Results}
\label{sec:results}

\subsection{Anatomy of One Execution}
\label{sec:anatomy}

Fig.~\ref{fig:trace} makes the object of evaluation concrete before any aggregate is reported. It retraces one complete execution---TEP IDV4, the fault whose evidence chain Section~\ref{sec:casestudy} examines in detail---exactly as the released run directory records it: the fingerprint written at ingestion, the ontology's variable semantics and its declared unmeasured channels, the statistical digest that found the recording's one large sustained shift---a cooling-valve step, the six hypotheses with four quantitatively eliminated alternatives, the parallel reviews, the endorsed report, and the finalization gate that turns the run into a gradeable result. Three properties of this trace carry the rest of the paper. Every decision-bearing output is a \emph{file}, not a conversational turn: the diagnostician argues from the statistical digest, the judge and the pre-auditor recompute their spot checks from the same artifacts, and no stage can rewrite another's output. The gates bind each transition: ontology before statistics, statistics before hypotheses, two independent reviews before the report is endorsed, and a mechanical inventory before the run counts as executed at all. And the trace is ordinary: this run cleared its diagnosis gate in the first pass and needed no repair round---true of twelve of the thirteen released runs, the exception being the judge-driven repair on TEP IDV3 (Section~\ref{sec:setup}). The remainder of this section aggregates the twelve executions this procedure produced.

\begin{figure}[!t]
\centering
\includegraphics[width=0.94\textwidth]{figures/fig_trace.png}
\caption{Anatomy of one complete execution (TEP IDV4, the reactor-cooling-step scenario of Section~\ref{sec:casestudy}), retraced from the released run directory. Each row is one pipeline stage: what the stage did, the decision-bearing content of its artifact (abbreviated; the exact strings are in the released artifacts), and the checkpoint its transition must satisfy (Table~\ref{tab:checkpoints}). Blue chips are agent stages, grey chips deterministic scripts; Steps 5a/5b executed as two concurrent sessions. The bounded repair edge (judge $<90$, dashed in Fig.~\ref{fig:pipeline}) was not triggered in this run. The Step-9 \textsc{pass} is the execution proof that contract C4 requires before a scenario may be graded (Section~\ref{sec:framework}).}
\label{fig:trace}
\end{figure}

\subsection{Overall Performance}
\label{sec:overall}
Table~\ref{tab:overall} reports the headline numbers. Executing the full pipeline under truth isolation (one recorded execution per scenario), IDD resolves \textbf{4/9} fault scenarios to a single root cause (Top-1 $=44.4\%$, Wilson 95\% CI $[18.9\%,73.3\%]$); in the remaining five, elimination did not reduce the surviving set to a single mechanism and no channel separates the survivors, so the pipeline returns capped \mbox{\textsc{competing\_set}} verdicts---\emph{in every one of the nine scenarios the true mechanism is named within the primary finding or the surviving-hypothesis set} (top-$k$ $=9/9$). Read at the keyword instrument's family level, the five failures are failures of \emph{resolution}---the evidence cannot separate the survivors---rather than omissions of the true mechanism family. The statement is set-membership, and its specificity is not tested by the instrument: the protocol mandates at least three hypotheses per scenario and the keyword sets derive from the same documented cause table (Section~\ref{sec:setup}), so family-level membership is close to design-guaranteed; top-$k$ is accordingly read as an instrument reading, never as a resolution bound (Section~\ref{sec:discussion}).

The keyword instrument needs three disclosures here, because one of the four resolved verdicts sits close to its grading tokens. The SKAB valve scenario's keyword set names the valve in a recording whose loop definitionally contains a valve; excluding it and batch 93, whose set contains the generic token ``deviation'', leaves Top-1 at 3/7 over the remainder, so the resolution pattern does not depend on the near-tautological tokens (top-$k$ on the same clean subset is 7/7); under a branch-localized top-$k$ reading, IDV3 counts as a miss and top-$k$ is 8/9. Batch 93 deserves the fuller statement: its verdict is genuinely execution-side and hits the truth set on \emph{actuator}---plus \emph{heating} and \emph{deviation}---but the diagnosis pins the cooling-water actuator at minimum flow while the documented deviation names the acid-dosing and heating-water flows, so under an identity-strict reading that verdict does not score and Top-1 over the nine faults is 3/9 (Wilson CI $[12.1\%,64.6\%]$). The batch-93 discount does not propagate to top-$k$ at the family level, which stays 9/9: set membership is exactly the mechanism-family instrument, and batch 93's execution-side family match is what that instrument measures. The branch-localized reading of Section~\ref{sec:setup} is stricter on the TEP side and does propagate: it discounts IDV3's stripper-branch membership, giving top-$k$ 8/9. Throughout, a hit certifies mechanism-family agreement, never mechanism identity. Under leave-one-token-out re-scoring of the released keyword-hit records, Top-1 moves to 3/9 only if the token ``header''---the single token on which IDV7's hit rests---is removed from that scenario's grading set; the reading is first-order, and multi-token removals only lower it further (IDV1's two hit tokens are spelling variants of one documented label): Top-1 is an instrument-dependent estimate bounded by 3/9--4/9 across the readings this section reports, and the token sets are released so each reading is reproducible. Because the keyword sets were calibrated on v1-era executions of these same scenarios (Section~\ref{sec:setup}), the family-hit readings are not an independent re-verification of the committed verdicts; the committed neutralized-identifier re-run will additionally re-check keyword coverage.

Determined-verdict precision is 4/4: no verdict is at once \mbox{\textsc{determined}} \emph{and} wrong. Under the identity-strict lens the batch-93 discount propagates here too: precision becomes 3/4 and overconfidence at least one---the asymmetry is reported rather than quoted in each metric's most favorable lens. All 3 control scenarios pass with \textbf{zero false alarms} (two of the three under briefs that disclose the segment's normal status; Section~\ref{sec:setup}). All 12 pipeline executions pass the finalization gate with valid event logs. The deterministic rubric averages \textbf{94.6/100} (96.2 excluding the mis-specified control check discussed above), with nine scenarios at 100; every deducted point traces to the disclosed instrument artifacts above, none to evidence substance. The deductions are SKAB inlet-valve throttling $-15$ (R4 verdict-type compliance, analyzed in Section~\ref{sec:calibration}; the 8/9 compliance count is itself an upper bound under the protocol's one-directional typing, which flags over-capping but cannot detect under-capping), TEP IDV1 $-15$, and the IndPenSim control $-35$. The two calibration deductions arise because reported confidence 0.91 and 0.93 sit $0.01$ and $0.03$ above the machine gate's flat 0.90 non-CS ceiling: the agent-side protocol documents no \mbox{\textsc{det}} cap, so the corpus exhibits the gap between the two cap sets that Section~\ref{sec:method} predicts. The control additionally loses 20 points on R3 grounding, its surviving ``normal operation'' hypothesis not being grounded in the fault-brief numbers the check targets. The in-pipeline judge gate averages 94.7 (range 90--98; all scores are post-threshold by the repair gate's construction), and the artifact-consistency gate verifies that the evaluation is internally consistent. Fig.~\ref{fig:benchmark} summarizes the outcomes; Section~\ref{sec:ablation} places them against model-controlled baselines.

\begin{table}[!t]
\caption{Overall Results on the 12-Scenario Benchmark (full-pipeline execution; one recorded execution per scenario, plus the random-draw re-execution of Section~\ref{sec:stability}). Rubric deductions, the two cap-set exceedances, and the identity-strict reading are analyzed in Section~\ref{sec:calibration} and Section~\ref{sec:overall}. The precision row conditions on the pipeline's own verdict choice; the printed interval corresponds to the identity-inclusive 4/4 reading, not to the identity-strict 3/4 shown alongside it. Judge scores are post-threshold ($\ge$90), averaged over $n{=}12$ graded scenarios, and carry no benchmark weight. Wilson 95\% intervals at $z{=}1.96$, matching the released scorer. Intervals are corpus-level measurement uncertainty: the nine fault scenarios are purposive and non-exchangeable (Section~\ref{sec:discussion}), and the grading keywords were calibrated on v1-era executions of these same scenarios (Section~\ref{sec:setup}). The top-$k$ row counts scenarios whose surviving-hypothesis set or primary finding names the true mechanism. The baseline arms of Section~\ref{sec:ablation} are likewise single recorded executions; ``strict'' for those arms counts any keyword hit, whereas the Top-1 row counts resolved verdicts only. Intervals are second-order: the dominant uncertainty is the reading structure itself (upper-bound, identity-strict, branch-localized), not sampling; the ceiling-compliance and judge-gate rows are gate properties enforced by construction, not truth-comparison outcomes.}
\label{tab:overall}
\centering
\footnotesize
\resizebox{\textwidth}{!}{%
\begin{tabular}{@{}p{6.7cm}p{2.3cm}p{2.3cm}@{}}
\toprule
Metric & Value & 95\% CI \\
\midrule
Resolved Top-1 (\mbox{\textsc{determined}} and correct; $n{=}9$) & 4/9 $=$ 44.4\% & [18.9, 73.3] \\
\quad Recall-possible bound on Top-1 (\S\ref{sec:discussion}) & upper 4/9, lower 0/9 &---\\
\midrule
\multicolumn{3}{@{}l}{\emph{Instrument readings of resolved Top-1 (Section~\ref{sec:overall})}}\\
Identity-strict reading (batch-93 verdict discounted) & 3/9 $=$ 33.3\% & [12.1, 64.6] \\
\quad Batch-93 excluded alone (truth anchor under-determined, \S\ref{sec:setup}) & 3/8 &---\\
Leave-one-token-out re-scoring (Section~\ref{sec:overall}) & 3/9--4/9 &---\\
\midrule
True cause within top candidates (top-$k$; $n{=}9$, design-aided) & 9/9 $=$ 100\% & [70.1, 100.0] \\
\quad Branch-localized top-$k$ reading (IDV3 discounted) & 8/9 &---\\
Determined-verdict precision (correct when \mbox{\textsc{determined}}) & 4/4 (identity-strict: 3/4) & [51.0, 100] cond.\ on the determined set \\
Ceiling compliance, capped-verdict states (confidence $\le$ CS cap) & 5/5 (gate-enforced by construction) &---\\
Ceiling compliance, non-CS states vs machine 0.90 gate & 5/7 (two exceed by 0.01/0.03; Section~\ref{sec:method}) &---\\
Verdict-type compliance (type in allowed set) & 8/9 & [56.5, 98.0] \\
Capped verdicts with the true mechanism ranked & 5/5 (five of the nine faults capped) &---\\
Control pass (no false alarm; $n{=}3$, one blind) & 3/3 & [43.9, 100.0] \\
Protocol-compliance rubric (author-designed checklist) & 94.6/100 mean & nine at 100, two at 85, one at 65 \\
Judge gate mean (in-pipeline, post-threshold $\ge$90) & 94.7 & range 90--98 \\
Pipeline finalization gate & 12/12 PASS &---\\
Artifact-consistency gate & verified &---\\
\bottomrule
\end{tabular}}%
\end{table}

\begin{figure}[!t]
\centering
\includegraphics[width=\textwidth]{figures/fig_benchmark.png}
\caption{Benchmark outcomes under full-pipeline execution. (a) Fault outcomes by dataset: Top-1-correct \mbox{\textsc{determined}} verdicts (dark blue) and capped \mbox{\textsc{competing\_set}} verdicts whose survivors rank the true mechanism (hatched sky blue)---no scenario's survivors exclude the true mechanism. Bar length is proportional to the dataset's fault count (the all-faults row, nine faults, is full scale); the printed fractions are keyword-instrument readings (Table~\ref{tab:overall} carries the identity-strict and recall-possible readings), and all three controls pass with zero false alarms (one of three blind; excluded from the bars). (b) Per-scenario deterministic rubric (filled markers) against the in-pipeline judge gate (open markers); dark blue = fault scenarios, grey = controls; the dotted line is the judge's 90-point gate. All 12 executions pass the finalization gate, and the artifact-consistency gate confirms that the evaluation is internally consistent.}
\label{fig:benchmark}
\end{figure}

\subsection{Per-Mechanism Breakdown and Within-Family Discrimination}
\label{sec:permech}
Table~\ref{tab:mech} lists the nine fault scenarios with the verdict each pipeline execution actually produced, the confidence, the rubric score, and the decisive statistical signature the diagnosis rested on. Three observations matter.

\textbf{Within-family form discrimination.} IDV4 (step), IDV11 (random variation) and IDV14 (valve stiction) all perturb the same physical channel---reactor cooling water---yet produce three distinct statistical forms that the pipeline separated, and resolved differently for recorded reasons. IDV4 (step) yields a cooling-triad step fingerprint (reactor temperature spike $z{=}{+}10.5$ with a sustained $+6.8\sigma$ valve step and a flat loop-external channel) whose sub-attribution, unmeasured inlet-temperature step versus sudden exchanger-UA loss (UA: the exchanger's overall heat-transfer conductance), the data cannot separate; the verdict is a capped \mbox{\textsc{competing\_set}}. IDV11 (random variation) yields a pure variance-burst signature with \emph{frozen means} (valve-position variance ratio $40.1$, temperature $16.4$, $\Delta$mean $\le 0.01\sigma$, actuator leading temperature by five samples at $r{=}{-}0.64$ while the closed-loop lag-0 sign flip is barred from causal citation) and is capped again on the same inlet-versus-fouling fork---the same conservatism as IDV4's: the frozen-mean burst disfavors every thermal-side account, but the elimination threshold (machine floor: $\ge$0.90 exclusion confidence with cited contradictions, and the eliminations this corpus records are all at $\ge$95\%) is not met for the UA-redistribution sub-case, so the fork is capped rather than forced, uniformly across the family. IDV14 (valve stiction) yields a self-excited stiction limit cycle (variance explosions of $\times$189/116/222 on three synchronised channels with \emph{frozen means}, quantised valve stepping at $633/800$ samples, and a lag structure that tightens from one to five steps) and resolves to \mbox{\textsc{determined}}, because the quantised-stepping plus limit-cycle signature---the classical stiction fingerprint of slip-jump and oscillation cross-correlation~\cite{horch1999simple,hagglund1995control,choudhury2008diagnosis}---is an actuator behavior no thermal hypothesis reproduces. Recalling the textbook label ``cooling-water fault'' alone therefore cannot produce these three different mechanism pictures; what the three pictures rule out is a \emph{channel-label} recall shortcut. The documented cause table does record each fault's dynamic form (Section~\ref{sec:discussion}), so the reading that carries evidential weight is not label recall but the one matched to, and argued from, the per-recording statistics that discriminate each form; the two capped verdicts name the one channel (cooling-water inlet temperature) whose addition would collapse their forks.

\textbf{Reasoning beyond the detection limit.} IDV3 and IDV4 lie beyond classical PCA monitoring and are excluded by FaultExplainer's PCA implementation~\cite{russell2000delay,chiang2001fault,khan2024faultexplainer}. IDV3---commonly regarded as the hardest TEP fault~\cite{russell2000delay}, with a deliberately weak signature (all column mean shifts $\le 0.92\sigma$)---terminates in a capped \mbox{\textsc{competing\_set}} (0.55): the exclusion chain eliminates the reactor block (flat: settling $z{=}{-}0.001$), the feed/load family, and sensor artifacts, while variance inflation concentrates in the stripper steam loop (variance ratios $2.1$--$2.7$, Levene $p{<}10^{-4}$, monotone stripper-temperature drift); the fork between a steam-side thermal disturbance and a stripper-feed thermal-state change cannot be resolved without unmeasured steam-header and feed-temperature measurements, and the verdict names them. The true mechanism family remains ranked among the survivors. IDV4 terminates similarly (0.63) with the heat-removal-step family top-ranked on an energy-balance necessity argument (Section~\ref{sec:casestudy}). Both caps fall where classical detection and the published baseline stop; the alternative, forcing a single cause from indistinguishable evidence, is the failure mode the protocol is designed to prevent.

\textbf{Blind-set robustness.} The SKAB cavitation scenario is drawn from the ``other'' experiment group whose fault label is withheld from the scenario description (mechanism-blind---all fault briefs assert that an anomaly exists, and this one withholds the mechanism word; not identifier-blind: the case id and run-directory name still carry it; Section~\ref{sec:discussion}). The pipeline's surviving hypotheses are the two sides of an intermittent flow-collapse instability (76 near-zero-flow rows in forty 1--3\,s bursts with the pump load essentially unchanged at $2.34$ vs $2.40$\,A, same-window vibration rises of $+5\%/+3\%$ on two accelerometers, and a $+0.23$\,K temperature rise at the pump (consistent with local dissipation---pump power or collapse heating---rather than a directional heat input); the drive side and the voltage confound are quantitatively excluded), but the channels that would separate valve-side throttling oscillation from suction-side cavitation or air entrainment---valve position and high-frequency suction pressure, the channels whose added value vibration- and acoustic-emission-based cavitation monitoring formalizes~\cite{alfayez2005detection}---are absent, so the verdict is a capped \mbox{\textsc{competing\_set}} (0.60) with the cavitation-type family ranked among the survivors (second of two, behind the valve-side hypothesis---a rank that carries no confidence meaning, Section~\ref{sec:casestudy}), naming the missing channels.

\begin{table}[!t]
\centering
\renewcommand{\baselinestretch}{1}\selectfont\scriptsize
\caption{Per-scenario results with the decisive statistical signature (abbreviated; full chains in the released run directories (\texttt{04\_diagnostics}) and rubric details in \texttt{rubric.json}). \textsc{det} = \mbox{\textsc{determined}}, \textsc{cs} = capped \mbox{\textsc{competing\_set}}; every \textsc{cs} verdict ranks the true mechanism among its survivors. XMV$n$ denotes TEP's $n$-th manipulated variable.}
\label{tab:mech}
\renewcommand{\arraystretch}{0.94}
\resizebox{\textwidth}{!}{%
\begin{tabular}{@{}lcc>{\raggedright\arraybackslash}p{4.5cm}@{}}
\toprule
Scenario (mechanism) & Verdict & Rubric & Decisive signature \\
\midrule
SKAB inlet-valve throttling & \textsc{cs} 0.60 & 85 & flow plateau $-3.13\%$ ($p{=}5.7{\times}10^{-158}$, rows 621--927); current co-drop $-0.115$\,A; pressure quantised \\
SKAB cavitation (blind set) & \textsc{cs} 0.60 & 100 & intermittent collapse (76 rows, 40 bursts of 1--3\,s); pump load held (2.34 vs 2.40\,A); twin vibration $+5\%/+3\%$ \\
TEP IDV1 A/C feed-ratio step & \textsc{det} 0.91 & 85 & A-feed valve/flow $+18.5\sigma$ lockstep ($r{=}0.9996$); C-side $-2.9/-6.9\sigma$; steam trails at lag 2 \\
TEP IDV3 D-feed temp.\ (hard) & \textsc{cs} 0.55 & 100 & mean shifts $\le 0.92\sigma$; stripper steam loop var $\times 2.1$--$2.7$ ($p{<}10^{-4}$); reactor flat; thermal fork unresolvable; truth family on stripper-feed branch \\
TEP IDV4 cooling (step) & \textsc{cs} 0.63 & 100 & cooling triad: temp $z{=}{+}10.5$, XMV\_10 $+6.8\sigma$ step, loop-external flat; energy balance forces inlet/UA fork \\
TEP IDV7 C-header pressure & \textsc{det} 0.78 & 100 & XMV\_4 $+25.1\%$ ($z{=}12.4$) with flow constant; gain $-20.1\%$ $\Rightarrow$ $\Delta$P ratio 0.64 \\
TEP IDV11 cooling (random) & \textsc{cs} 0.60 & 100 & variance burst VR $40.1/16.4$, means frozen ($\le 0.01\sigma$); actuator leads temp 5 samples ($r{=}{-}0.64$); lag-0 flip barred \\
TEP IDV14 valve stiction & \textsc{det} 0.80 & 100 & VR $\times$189/116/222, means frozen; quantised stepping 633/800; lag tightens 1$\to$5 ($r{=}0.998/{-}0.964$) \\
IndPenSim batch-93 & \textsc{det} 0.82 & 100 & CW loop pinned (2.0 vs 193--700\,L/h) 72/72 rows $\to$ 301.4\,K; release $\to$ in band in 2.4\,h \\
\midrule
Controls (3) & 3 \textsc{det} (normal) & 100/100/65 & normal verdicts at 0.86--0.93; zero false alarms; raw correlations refuted as control signatures \\
\bottomrule
\end{tabular}%
}
\end{table}
\subsection{Ceiling Compliance and False-Alarm Suppression}
\label{sec:calibration}

\begin{figure}[!t]
\centering
\includegraphics[width=\textwidth]{figures/fig_calibration.png}
\caption{Reported confidence per scenario in the full-pipeline benchmark run. Solid dark blue bars are \mbox{\textsc{determined}} fault verdicts; hatched sky blue bars are the five capped \mbox{\textsc{competing\_set}} verdicts (all with the true mechanism ranked among survivors); hatched grey bars are the three normal-operation controls; the dashed line is the 0.70 \mbox{\textsc{competing\_set}} ceiling and the dotted line the machine grading gate's flat 0.90 non-CS ceiling. The two bars crossing the dotted line (TEP IDV1 at 0.91, $+0.01$; IndPenSim control at 0.93, $+0.03$) mark the corpus's one visible gap between the agent-side cap set (which documents no \textsc{det} cap) and the machine gate (Section~\ref{sec:method}); the grading rubric records both as R4 calibration deductions. Each capped verdict names the discriminating measurement that would separate its surviving mechanisms; no conclusion is simultaneously \mbox{\textsc{determined}} and wrong (overconfidence 0).}
\label{fig:calib}
\end{figure}

\begin{figure}[!t]
\centering
\includegraphics[width=\textwidth,height=0.78\textheight,keepaspectratio]{figures/fig_casestudy.png}
\caption{Worked example of a capped verdict (SKAB inlet-valve throttling, the benchmark's single verdict-type compliance miss). Top: the recorded flow signal holds a plateau in the two lowest quantisation levels ($\approx$30--31 vs $\approx$32\,L/min; $-3.13\%$, $p{=}5.7{\times}10^{-158}$; pale-amber shading marks the rows 621--927 plateau window) while the pump current co-drops only directionally ($-0.115$\,A, $p{=}0.10$) and vibration stays flat; the pressure channel's 0.327-bar quantisation (annotated in the panel) masks the head-direction measurement that would separate the two surviving mechanisms. Bottom strip: the evidence chain as the pipeline recorded it, from the raw signal through L3 statistics and the L5 mechanism to the two surviving hypotheses (downstream throttling 0.65 $\|$ upstream supply restriction 0.60---ranks within a capped set carry no confidence meaning) and the capped verdict with the named missing channels. Every annotation is read from the released run artifacts.}
\label{fig:casestudy}
\end{figure}

Fig.~\ref{fig:calib} plots the reported confidence per scenario. Fault-scenario confidences span 0.55--0.91; the five \mbox{\textsc{competing\_set}} verdicts sit in the 0.55--0.63 band because the protocol caps them, and each names its discriminating measurement (valve-position/command channel and pressure direction; suction-side high-frequency pressure; cooling-water inlet temperature for both cooling-loop forks; steam-header and feed-temperature measurements for IDV3). Within the surviving sets the documented truth ranks first in two of the five capped verdicts at family or branch granularity (IDV4; IDV11 at branch level, where the cooling-water-side disturbance heads the set while its specific random-variation form ranks third of four behind the steam-branch alternative); at form-specific granularity the count is one of five. The truth ranks second in the remaining three (SKAB inlet-valve throttling and cavitation, each of two survivors; IDV3, of three, ahead of only a weak-excluded residual candidate). The cap-internal order has one defined meaning: it reflects the five-factor evidence weights at capping time. It is not a calibrated confidence, carries no truth guarantee, and top-$k$ counts set membership, never ordering. Per-survivor ranks and five-factor scores are recorded in the released gradings, so the orderings are independently reconstructable. We report this metric as \emph{ceiling compliance} rather than calibration; it is circular by construction, since the agent-side caps are enforced and compliance with them is a property of the protocol's arithmetic.

The non-CS states are where the corpus contributes new evidence: the agent-side protocol documents no \mbox{\textsc{determined}} cap, while the machine grading gate sets a flat 0.90 flag bound---and two verdicts exceed it, TEP IDV1 at 0.91 and the IndPenSim control at 0.93. Both are recorded as R4 calibration deductions by the rubric and plotted in Fig.~\ref{fig:calib}; neither verdict is wrong (the overconfidence metric---\mbox{\textsc{determined}} \emph{and} wrong---never fires), but the two-point gap between the agent-side and machine-side cap sets that Section~\ref{sec:method} predicted is now visible in the corpus rather than hypothetical---a consistency finding about our own two cap artifacts, not external validation. The committed remedy is structural: a single machine-readable cap table consumed by both the protocol writer and the grading gate, which removes the drift class instead of reconciling its instances. A psychometric calibration study (pre-gate confidence versus outcome correctness, over repeat runs) remains future work. Throughout this corpus, reported confidence is a \emph{protocol ordinal}: the capped five-factor assessment, not a probability of correctness and not a coverage statement---$0.60$ does not mean six verdicts in ten are right. The natural formalization of the set-valued \mbox{\textsc{competing\_set}} output---a distribution-free guarantee on the survivor set via conformal risk control~\cite{angelopoulos2023conformal}---is listed with the calibration study as future work.

Under the benchmark's strict typing, one scenario counts as a verdict-type miss: SKAB inlet-valve throttling returns \mbox{\textsc{competing\_set}} where the scenario's allowed set admits only \mbox{\textsc{determined}} (the ground truth is known to be valve restriction). The typing is one-directional: the other eight fault scenarios' allowed sets admit both states, so compliance can flag over-capping but never under-capping. The miss stands; the capped verdict is nonetheless the act the protocol demands when the valve-position and controller-command channels that would separate ``valve closing'' from ``commanded duty reduction'' are absent from the recording: the pump's electrical load shows only a directional co-drop with the flow plateau ($-0.115$\,A, $p{=}0.10$), the pressure channel is quantised at 0.327 bar (roughly an order of magnitude above the expected head change), and the flow level shift itself is small and unambiguous in timing but silent about \emph{which} side of the restriction moved. Fig.~\ref{fig:casestudy} makes this scenario fully visible: the annotated recording (flow plateau in the lowest quantisation levels with the pump load essentially flat) alongside the evidence chain as the pipeline recorded it---raw signal, statistics, mechanism, the two surviving hypotheses, and the named missing channels.

Control scenarios demonstrate the complementary discipline: the normal IndPenSim batch contains recipe-dosing impulses reaching $|z|{=}16.2$, which the pipeline attributes to recipe steps and closed-loop control actions rather than to faults; the normal TEP run exhibits feedback signatures at $|r|\ge 0.95$ with the loop's lag-0 peak correctly read as controller response; the normal SKAB run's strong raw correlations ($|r|$ 0.71--0.89) are refuted as slow co-drift artifacts by the detrending and lag-compensation checks. These are the failure modes that produce false alarms in naive correlation-driven diagnosis.

\subsection{Case Study: An Auditable Evidence Chain on a Fault Excluded from Classical Scoring}\label{sec:casestudy}
For TEP IDV4---a fault excluded from FaultExplainer's scoring by its PCA implementation, and only weakly detected classically ($T^2$ 48.8\% vs SPE 100\% in our deterministic replication)~\cite{russell2000delay}---the recorded chain reads: (i)~\emph{L3 statistical evidence}: at fault onset the reactor cooling valve XMV\_10 (XMV$n$: TEP's $n$-th manipulated variable, i.e.\ valve position) takes a sustained $+6.8\sigma$ step (41.1\%$\to$44.9\% with no ramp) while reactor temperature spikes to $z{=}{+}10.5$ with a transient on its sister thermocouple ($+2.0/-2.3\sigma$) and every loop-external column stays flat (the separator-outlet column at $-0.08\sigma$)---a loop-localized cooling triad; (ii)~\emph{L5 mechanism, with a necessity argument}: an energy balance on the reactor heat-removal loop shows the observed coolant-flow rise with conserved reactor temperature algebraically \emph{requires} an upward shift in the cooling-water inlet temperature (or, indistinguishably, a sudden drop in exchanger UA)---a controller compensating exactly as the physics demands, not a sensor artifact; (iii)~\emph{eliminations}: feed-side, condenser, purge, and ramp-form alternatives are quantitatively excluded (each at $\ge$95\% confidence with cited contradictions), and the feed-composition channel shows no in-window signature; (iv)~\emph{conclusion}: capped \mbox{\textsc{competing\_set}} at 0.63 (INDISTINGUISHABLE ceiling 0.65) between ``cooling-water inlet-temperature step'' (top-ranked) and ``sudden UA loss''. The recording contains no cooling-water inlet measurement; the verdict names it, and the audit trail quantifies the fork: conserving the loop temperatures under the UA sub-case requires a pre-existing reactor-to-coolant temperature approach below $1.5$\,K (implying a coolant inlet above $\approx$93\,\textdegree C, against the $\approx$26\,K approach observed in the pre-onset data), which strongly disfavors the pure-UA form---though it stays below the quantitative elimination threshold because a UA loss can redistribute the loop's LMTD rather than move the observed temperatures, which is precisely why the verdict is a capped competing set rather than a determination. The quantification sharpens the ranking within the cap. The full chain, together with the physics-check verdict, the judge's ten-criterion scores (97/100), and the independent HTML review (pass), is persisted in the grading record and released with the evaluation. We read the cap as the finding: on a fault the published baseline cannot even score, the pipeline localizes the loop, states the physically necessary mechanism family, quantifies why the two remaining sub-mechanisms cannot be separated, and names the single measurement---the cooling-water inlet thermocouple---whose installation would collapse the set to \mbox{\textsc{determined}}.

\section{Comparison With Published Baselines, Model-Controlled Ablation, and Verdict Stability}
\label{sec:comparison}

Section~\ref{sec:results} reported what the pipeline concluded on its own recorded terms; this section subjects those conclusions to progressively harsher external checks---the published per-fault record, the model-controlled baseline suite that bounds every accuracy claim (Section~\ref{sec:ablation}), and a drawn re-execution that probes verdict stability (Section~\ref{sec:stability}).

\subsection{Register A: Per-Fault Comparison on TEP}
FaultExplainer~\cite{khan2024faultexplainer} reports, in its root-causes-included regime, root-cause correctness of 7/11 (63.6\%, Wilson CI $[35.4\%,84.8\%]$) for GPT-4o and 9/11 (81.8\%, CI $[52.3\%,94.9\%]$) for o1-preview on the 11 PCA-detectable TEP faults---correctness meaning the true fault or an alias in the model's top-3, with candidates supplied in the prompt; in its no-candidate regime both models reach 8/11 under lenient ``related'' grading. IDD, without candidate lists and scored by strict mechanism keywords, is Top-1 correct on \textbf{3/6} TEP faults (Wilson CI $[18.8\%,81.2\%]$) and ranks the true mechanism among the survivors in all six (top-$k$ 6/6); on the four shared PCA-detectable faults it is correct on 3/4, where GPT-4o is 3/4 (missing IDV1, which IDD resolves) and o1-preview is 4/4 under alias-inclusive scoring. The base models differ and are not controlled, the grading protocols differ (Section~\ref{sec:setup}), and the ablation of Section~\ref{sec:ablation} shows that a bare call to the same model family already scores these documented faults at or above the pipeline (9/9 versus 4/9 resolved)---so we read the per-fault table descriptively, not as a superiority claim. Table~\ref{tab:perfault} (\ref{app:artifacts}) transcribes the comparison; AgentRCA's TEP evaluation set is not per-fault alignable with ours, so it enters the positioning of Table~\ref{tab:related} rather than this comparison; Fig.~\ref{fig:perfault} breaks it down per fault against FaultExplainer's Table~1 outcomes and PCA detectability, and Fig.~\ref{fig:forest} shows the aggregates with intervals. Most of the comparison concentrates in three outcomes:

\begin{itemize}
\item \textbf{IDV1 (A/C feed-ratio step):} the GPT-4o baseline mis-diagnoses this fault (o1-preview resolves it); IDD resolves it to \mbox{\textsc{determined}} (0.91) from the actuator--flow signature (A-feed valve and flow in $+18.5\sigma$ lockstep, $r{=}0.9996$, with the C-side feed compensating at $-2.9/-6.9\sigma$), correctly restructuring the disturbance as an A:C supply-weight shift with the stripper steam response trailing as a downstream interlock (lag 2).
\item \textbf{IDV3 and IDV4 (excluded from FE scoring):} both faults are excluded from FaultExplainer's scoring by its PCA implementation. IDD terminates both in capped \mbox{\textsc{competing\_set}} verdicts with the true mechanism ranked (Section~\ref{sec:casestudy} details the IDV4 chain).
\item \textbf{IDV7, IDV11, IDV14:} IDD resolves the header-pressure form (0.78: valve--flow gain loss with flow constant $\Rightarrow$ supply-capacity drop) and the stiction form (0.80: quantised-stepping limit cycle), and caps IDV11's random-variation form (0.60) at the inlet-vs-fouling fork---separating the three cooling-family \emph{forms} at the level of distinct quantitative signatures (Section~\ref{sec:permech}), a granularity below fault-identity level.
\end{itemize}

Because the scenario sets, prompt regimes, and grading protocols all differ (top-3 with alias acceptance vs.\ a strict single verdict; with vs.\ without candidate lists; different base models), we present the per-fault table as descriptive evidence and claim no significance test across systems. The comparison is also asymmetric in cost: FE-style prompting runs in minutes per fault, whereas the full gated pipeline's per-scenario agent time is on the order of two hours (Section~\ref{sec:setup}).

\begin{figure}[!t]
\centering
\includegraphics[width=0.88\textwidth,height=0.68\textheight,keepaspectratio]{figures/fig_tep_perfault.png}
\caption{Per-fault comparison on the TEP subset: IDD vs.\ FaultExplainer (transcribed Table~1 outcomes of~\cite{khan2024faultexplainer}; top-3 with alias acceptance) and PCA detectability as used in FaultExplainer's scoring (Russell et al.~\cite{russell2000delay}; Chiang et al.~\cite{chiang2001fault})---the column reproduces FE's $T^2$-only screen as documented in its release, not detectability in our own $T^2$+SPE re-analysis. IDD is Top-1 correct on 3/6 without candidate lists (IDV1, IDV7, IDV14) and ranks the true mechanism among the survivors in all six; the three capped \mbox{\textsc{competing\_set}} verdicts name their discriminating measurements. Our own re-analysis does detect IDV4 ($T^2$ 48.8\%, SPE 100\%) while IDV3 remains at the detection limit (2.9\%/4.8\%); because scenario sets and grading protocols differ, no significance test is claimed across systems. Cell colors: solid dark blue = IDD Top-1-correct; hatched sky blue = IDD capped \mbox{\textsc{competing\_set}} (true mechanism ranked); green = FE-correct; vermilion = FE-wrong; grey = FE-unscored or PCA-excluded.}
\label{fig:perfault}
\end{figure}

\begin{figure}[!t]
\centering
\includegraphics[width=\textwidth]{figures/fig_forest.png}
\caption{Register-A aggregate comparison with Wilson 95\% intervals ($z{=}1.96$). Grey: FaultExplainer-reported single-LLM baselines ($n{=}11$, candidates in prompt). Blue: IDD full-pipeline execution without candidate lists---the TEP subset (3/6), all nine faults Top-1 (4/9), and top-$k$ (9/9, every fault's survivors rank the true mechanism). Green: the three controls (3/3, no false alarms). Open markers mark the two readings bounded at 100\% (top-$k$, controls), whose upper interval limits sit at the axis edge. Sets differ, so overlapping intervals are expected and no significance test is claimed; scenarios are purposive and non-exchangeable, so intervals are corpus-level measurement uncertainty (Section~\ref{sec:discussion}); the per-fault breakdown (Fig.~\ref{fig:perfault}; Table~\ref{tab:perfault}) is the finer-grained evidence.}
\label{fig:forest}
\end{figure}

\subsection{Register B: Reference-Only Context}
For context we report sample-level results from the literature without ranking: TEP classification accuracies of 0.887--0.907 (MLP/GRU/TCN on the 28-class extended TEP simulation of~\cite{rieth2017additional}, as evaluated in~\cite{pozdnyakov2024adversarial}), anomaly-detection F1 of 0.970 (BeatGAN~\cite{zhou2019beatgan}, as evaluated in~\cite{hartung2023deep}), SKAB ensemble F1/AUC of 0.85/0.88~\cite{iliopoulos2023detection}, and leakage-controlled bearing AUROC of 0.75--0.93~\cite{vieira2026towards}. These systems answer ``is this sample anomalous / which class''; IDD answers ``what is the mechanism, how confident, with which evidence''. The dimensions are complementary: a deployment would plausibly chain a sample-level detector (first-stage alarm) with a scenario-level RCA system (this work).

\subsection{Model-Controlled Baselines and Pipeline Ablation}\label{sec:ablation}
Because the base model is a confound in any comparison with published LLM baselines (Section~\ref{sec:setup}), we ran three model-controlled baselines with the \emph{same model family and deployment} as the pipeline (Table~\ref{tab:ablation}; the LLM arms are recorded replays of live runs from this deployment, and the TEP fault-free control was not run bare): (i) a classical PCA baseline---normal-control training, $T^2$/Q detection with 99th-percentile alarms, and SPE-contribution diagnostics~\cite{chiang2001fault,qin2012survey}---on all twelve scenarios; (ii) a replication of FaultExplainer's protocol on the six TEP faults, in both of FE's prompt regimes, scored both FE-style (top-3 with alias classes) and strictly; and (iii) the pre-registered single-LLM ablation---one call to the same model family over the identical blinded brief, whose evidence block is the pipeline's deterministic statistical digest (per-column anomaly profile and top correlations; no ontology, no pipeline stages, no gates)---on all nine fault scenarios, scored with the same strict mechanism keywords as IDD. All three baselines are bare single-call executions: no agent orchestration, no gates, no evidence grading---and, like the pipeline runs, each is a single recorded execution rather than a repeated sample. The classical PCA arm is computed by a released deterministic script (pure-Jacobi eigendecomposition, normal-control training, 99th-percentile alarms of the control distribution, TEP fault-window evaluation from row 161); an earlier prototype run agreed with it on every SPE statistic and control rate, with minor $T^2$ differences on the weakest-signature faults, and the prototype output is archived in the released repository for audit.

\begin{table}[!t]
\caption{Model-Controlled Baselines on the Benchmark (same model family and deployment as the IDD pipeline; single call per scenario; answers in the released repository). The bare-LLM and FE-prompt arms are recorded replays of live runs; their feature extraction and prompt assembly, like the classical PCA arm, are deterministic re-executions. Fig.~\ref{fig:matrix} shows the same arms per scenario; this table aggregates by arm. Wall-clock: the single-call arms complete in minutes per scenario, against the pipeline's $\approx$1.9\,h median agent span.}
\label{tab:ablation}
\centering
\footnotesize\setlength{\tabcolsep}{3pt}
\resizebox{\textwidth}{!}{%
\begin{tabular}{@{}p{3.4cm}cp{2.3cm}c>{\raggedright\arraybackslash}p{3.2cm}@{}}
\toprule
Baseline & Candidates & Scoring & Scope & Result \\
\midrule
Classical PCA (normal-control training) & --- & detection rate & 12 scen. & IDV3 $T^2$ 2.9\%, \mbox{SPE 4.8\%}; IDV4 \mbox{SPE 100\%} vs $T^2$ 48.8\%; controls 1\% \\
Classical PCA (diagnosis) & --- & SPE contribution & 12 scen. & top-3 variables only; no mechanism verdict \\
FE-protocol replication (GLM) & yes & FE-style top-3 + alias & TEP 6 faults & 6/6 \\
FE-protocol replication (GLM) & yes & strict single-verdict & TEP 6 faults & 6/6 \\
FE-protocol replication (GLM) & no & strict single-verdict & TEP 6 faults & 6/6 \\
FE official prompts (GLM), T$^2$-contribution features & yes & FE-style top-3 + alias & TEP 6 faults & 5/6 \\
FE official prompts (GLM), T$^2$-contribution features & yes & strict single-verdict & TEP 6 faults & 5/6 \\
Single-LLM, identical blind digest (GLM) & no & strict single-verdict & 9 faults & 9/9 \\
Single-LLM, controls (GLM) & no & normal verdict & 2 controls & 2/2, zero false alarms \\
\textbf{IDD full pipeline} & \textbf{no} & \textbf{strict resolved Top-1} & \textbf{9 faults} & \textbf{4/9} (5 capped \mbox{\textsc{competing\_set}}, true mechanism ranked in all) \\
\bottomrule
\end{tabular}%
}
\end{table}

Three results follow, and they do not all favor the pipeline. \textbf{First}, on keyword-scored accuracy over documented public faults the bare same-model call is not merely equal to the full pipeline but ahead of it: 9/9 strict hits versus the pipeline's 4/9 resolved (five capped verdicts; on the same token-excluded clean subset as Section~\ref{sec:overall}'s Top-1 reading, the bare arm is 7/7)---including IDV3, plausibly by literature recall of the documented weak-signature fault, and naming valve and cavitation mechanisms on the SKAB scenarios from their signatures---and the recall caveat cuts both ways, since the pipeline's keyword sets derive from the same documented cause table: neither arm's absolute TEP numbers measure contamination-free ability, and the form-discrimination evidence of Section~\ref{sec:permech} together with the novel-disturbance suite is the clean test. The bare arm produces no false alarms on the two of the three controls the bare-LLM arm covered (the TEP fault-free control was not run bare; the released baseline record notes the coverage gap, and only covered controls are scored). The gap between 9/9 and 4/9 is the paper's central cost statement: the pipeline gives up resolution to gain an auditable process, and the capped verdicts are where that cost is visible. The cost is computable from the released gradings without new runs: had the protocol forced commitment to each capped verdict's top-ranked survivor, Top-1 would be 6/9 with three wrong commitments (both SKAB capped cases and IDV3 rank the truth second)---the cap trades three false determinations for two foregone resolutions, which is the trade the three-state protocol is built to make. In selective-prediction terms, the bare call and the forced-commitment policy are two operating points on one risk--coverage curve, and the capped protocol deliberately selects the conservative point \cite{chow1970optimum,angelopoulos2023conformal}. (The count is robust to that ordering's caveat: committing uniformly at random within each capped set yields an expected $\approx$6/9---the recorded orderings convert the expectation into 6/9 with three specific wrong commitments; they do not manufacture it.) That the resolution loss locates in the cap-and-elimination policy rather than in orchestration is inferred from this counterfactual, not tested by an arm; the filter-disabled arm remains pre-registered.

\textbf{Second}, the FE official-prompt arm, running on the same GLM model, scores 5/6 on the TEP subset---above IDD's 3/6 Top-1---with a different miss pattern: it misses IDV4 because its PCA $T^2$ features are dominated by feed-composition shifts that mask the cooling-water signal, while IDD caps IDV3, IDV4 and IDV11 (the two cooling-loop forks and the steam-site fork), keeping the true mechanism ranked in each. Read at IDD's closest analogue to FE's top-3 notion---the surviving-hypothesis set---the pipeline's top-$k$ on the TEP subset is 6/6.

\textbf{Third}, the classical PCA baseline reproduces the literature (IDV3 at the detection limit, $T^2$ 2.9\%/SPE 4.8\%; IDV4 ``weak $T^2$, strong $Q$''---SPE 100\% vs $T^2$ 48.8\%) and produces no mechanism verdicts.

Table~\ref{tab:baselinedetail} (\ref{app:artifacts}) opens both baseline arms to per-scenario inspection, beyond the hit/miss matrix of Fig.~\ref{fig:matrix}: it lists the bare call's actual top-1 answer per scenario and the PCA arm's SPE-contribution top-3 variables, alongside the IDD verdict. Three contrasts deserve a close reading. \emph{First}, the bare call commits one unqualified answer per scenario---no confidence, no verdict type, no evidence chain, no named missing measurement---and on SKAB cavitation it ranks cavitation \emph{first}, whereas IDD caps the same scenario with cavitation ranked second of two behind the statistically favored valve-side hypothesis (Section~\ref{sec:permech}): the same mechanism appears, but only the capped form states that the recording cannot separate the two and names the channels that would. \emph{Second}, on IndPenSim batch-93 the bare call names the acid-feed pH loop and the PCA contribution list is headed by the heating-water flow---both nearer the documented deviation description than IDD's cooling-loop reading, which is precisely why IDD's verdict discounts under the identity-strict reading; the arms' divergence on this scenario is the ambiguity the three-state protocol is built to surface rather than paper over. \emph{Third}, the PCA contribution lists are variable-level only, and variable ranking is not mechanism attribution: the implicated actuator XMV\_10 surfaces for IDV4 and IDV11; for IDV14 the list is headed by the reactor cooling-water outlet temperature (XMEAS\_21)---the affected loop---while the implicated valve XMV\_14 itself does not appear in the top-3 at all, the detection--diagnosis gap in its most concrete form; for the SKAB valve case the flow channel ranks third behind two vibration channels; and no list carries a confidence, a verdict state, or an exclusion chain. The per-scenario answers are released verbatim (one is in Chinese; the table uses abbreviated English glosses), so each of these readings is checkable against the record without re-running anything.

The three failures differ in kind, and the difference exposes what benchmark accuracy on documented faults conflates: detection ability, mechanism attribution, and epistemic discipline. The bare call produces nothing the audit plane can verify: no evidence levels, no ceiling compliance, no physics check, no execution proof, no artifact provenance---and its IDV3 answer, being correct, is exactly the kind of unverifiable guess that the three-state protocol is meant to constrain (the pipeline's capped verdict names the steam-header and feed-temperature measurements that would settle the question).

\begin{figure}[!t]
\centering
\includegraphics[width=\textwidth]{figures/fig_suite_matrix.png}
\caption{Per-scenario comparison matrix across the model-controlled baseline suite (all arms executed on the same prepared recordings; every cell read from the released run artifacts). Column 1: the pipeline's verdict with reported confidence---dark blue resolved \mbox{\textsc{determined}}, hatched sky blue capped \mbox{\textsc{competing\_set}} with the true mechanism ranked (all five capped verdicts carry that property), green ``normal'' controls. Column 2: bare single-call LLM under strict mechanism keywords (``--'' on the TEP control = coverage gap, not a pass). Column 3: the same model under FaultExplainer's candidate-list regime, FE-style scoring. Columns 4--5: classical PCA $T^2$/SPE detection rates from the suite's live re-execution on the same prepared recordings (white-to-green shading scales with the rate; controls sit at the nominal $\approx$1\% false-alarm level). Column 6: FaultExplainer's six-consecutive-$T^2$ trigger (the green ``no'' on the fault-free control is the correct outcome; ``--'' = arm not run outside TEP)---a detection result, distinct from the scored FE-official answers of Table~\ref{tab:ablation}.}
\label{fig:matrix}
\end{figure}

The classical PCA baseline quantifies the detection--diagnosis gap directly: IDV3 sits at the detection limit (2.9\% $T^2$, 4.8\% SPE), IDV4 is detected reliably only by the SPE statistic (100\% vs 48.8\%), and the output stops at variable contributions---no mechanism, no confidence, no verdict state. Attribution of the pipeline's remaining margins (e.g., whether the anti-spurious filter or the ontology layer contributes) still requires the filter-disabled and engine-substitution arms, which remain future work. A reproducibility note completes the arm: after the benchmark closed, the deterministic suite arms (classical PCA, FE-protocol feature extraction, prompt assembly) were re-executed end to end---a full consecutive double run of the 45-run suite---and reproduced all 45 run outputs byte-for-byte with timestamps excluded (released verdict: DETERMINISTIC, 45/45 identical; the re-execution covers the deterministic suite arms, not the stochastic reasoning layer), so the baseline numbers are exactly as reproducible as the deterministic layer they run on. Fig.~\ref{fig:matrix} arranges the complete comparison as a per-scenario matrix: the pipeline's verdict column against the two LLM-baseline arms and the three classical-detection statistics, making visible both where the systems agree (all nine faults are keyword-hit by the bare LLM, and the covered controls pass---two of three; released baseline record) and where they differ in kind (the pipeline's five capped verdicts, all with the true mechanism ranked, and the classical detectors' blindness to IDV3 and to IDV4's $T^2$ channel).

\subsection{Verdict Stability: a Random-Draw Independent Re-execution}\label{sec:stability}
The evaluation framework formalizes consistency testing as a draw-then-re-execute protocol: a scenario is selected by a \emph{seeded uniform random draw} over the nine-fault pool (so the selection is reproducible and fixed before comparison, though the seed originates from the evaluation harness rather than an external randomness beacon---a limitation we state alongside the draw record) and then re-executed through the full gated pipeline in a fresh run directory with fresh agent sessions (new ontology build, new statistical stage, new diagnosis, new review cycle, new HTML build and review); the draw record is append-only and replayable from the released seed. In the current corpus the draw selected TEP IDV11 (seed \texttt{1003818694}; the record carries the pool, the raw uniform value, and the resulting index), and the drawn re-execution completed all nine stages end to end in $\approx$2.0\,h, passing every gate (finalization \textsc{pass}). The framework's rule is that the latest proven execution of a scenario is its canonical graded run, so the re-execution (0.60) superseded the original batch execution (0.62) in every aggregate; both are capped \mbox{\textsc{competing\_set}} verdicts with the same signature, so the supersession changes no headline number. Fig.~\ref{fig:consistency} maps the full audit.

The two executions of the drawn scenario---the corpus's single repeat pair ($n{=}1$): the original batch execution \texttt{202609171639527} (run~1) and the drawn re-execution \texttt{202609172247007} (run~2, now canonical)---agree on everything the protocol measures at the mechanism level: the verdict state (capped \mbox{\textsc{competing\_set}} in both), the implicated actuator (XMV\_10, the reactor cooling-water valve), the mechanism class (environment-side cooling disturbance), and the top-ranked hypothesis (a cooling-water-side disturbance rather than a valve-internal fault), with the two sessions' confidence estimates 0.02 apart (0.62 vs 0.60, both under the 0.65 INDISTINGUISHABLE cap). Both sessions also agree on the discriminating evidence: the variance-burst topology with frozen means and the actuator-leads-temperature lead structure; the retest additionally tightened the causal-direction argument (lag-5 $r{=}{-}0.636$ with the closed-loop lag-0 sign flip formally barred from causal citation) and bounded the excluded feed and condenser families more sharply. The consistency audit marks the pair \textsc{consistent}---the current corpus holds exactly one case-version pair, and within-era agreement is therefore 1/1 with zero divergences; future repeats fold $\Delta$confidence and evidence-chain depth into the criterion, per the pre-registered repeat metrics of Section~\ref{sec:discussion}. The agreement is decided at the structured level (verdict type, equipment tag, mechanism class); the two sessions' free-text mechanism strings share only 0.14 of cause tokens, and the re-run's specific form ranks within the same cooling branch---wording drifts, the structure does not. One consistent pair is the beginning of a stability study, not its end. The agreement rule evaluates the coarsest structured fields (verdict type, equipment tag, mechanism class), and the drawn scenario is one whose deterministic signature---a variance burst of ratio 40 with frozen means---makes structural agreement close to mechanical, since both sessions run the same statistics stage; a draw landing on IDV3 or a SKAB case would have stressed the rule harder. With one draw, the protocol's evidential yield is itself unmeasured. The supersession rule (latest proven execution is canonical) was fixed before comparison and will not be applied selectively to disfavor re-runs; the remaining eleven fault-or-control scenarios each hold a single current-era execution, the audit lists them as outstanding re-run contracts rather than assuming their stability, and the released consistency report enumerates them. Expanding the repeat count (all 12 scenarios, $k \ge 3$ repeats, verdict-state stability and confidence distributions) remains pre-registered in the released repository.

\begin{figure}[!t]
\centering
\includegraphics[width=\textwidth,height=0.68\textheight,keepaspectratio]{figures/fig_consistency.png}
\caption{Consistency audit of the current-era corpus, read from the released audit artifact. (a) Each marker is one proven execution at its reported confidence: dark blue filled circles \mbox{\textsc{determined}} fault verdicts, sky filled squares capped \mbox{\textsc{competing\_set}} fault verdicts, grey open diamonds controls (the asterisk marks the brief-blind control); TEP IDV11 carries two adjacent markers (canonical execution and drawn re-execution, $\Delta$confidence $= 0.02$). Dashed line: the 0.70 \mbox{\textsc{competing\_set}} cap. (b) The re-test pair: verdict state, implicated actuator (XMV\_10), and mechanism class match across independent sessions; audit verdict \textsc{consistent}. The draw was seeded and replayable (seed released), executed before any comparison. The eleven single-execution scenarios carry explicit re-run contracts rather than assumed stability.}
\label{fig:consistency}
\end{figure}

\section{Discussion, Limitations, and Deployment Prospects}
\label{sec:discussion}

This section consolidates boundaries that the results already carry, in four groups: evidence strength (single executions, purposive sampling, the reading structure of Section~\ref{sec:overall}), prompt-knowledge contamination, the operating semantics of capped verdicts, and the provenance, cost, and deployment boundary of the audit plane.

\textbf{Statistical power and single-execution evidence.} With $n{=}9$ fault scenarios, Top-1 $=44.4\%$ supports only a wide interval ($[18.9\%,73.3\%]$ at 95\% confidence); the top-$k$ bound (9/9, $[70.1\%,100\%]$) is the stronger statement the current sample size can support---read as family-level set membership, not as a resolution bound: the instrument's strictness is not specificity, since the protocol mandates at least three hypotheses and the keyword sets derive from the documented cause table (Section~\ref{sec:setup}); what 9/9 certifies is that no failure omitted the true family, not how often the protocol could have committed earlier. Scenario selection was purposive rather than random: the suite over-represents documented, public fault episodes, which biases difficulty downward and buys mechanism diversity across families rather than within them, so the intervals in Table~\ref{tab:overall} read as measurement uncertainty on this corpus, not sampling generalization to a scenario population. Three qualifications attach to these figures. First, every headline number is a \emph{single recorded execution} per scenario---except TEP IDV11, whose drawn re-execution (Section~\ref{sec:stability}) reproduced the verdict state and mechanism at $\Delta$confidence 0.02, the only repeat measurement the corpus currently holds. Second, the bounded repair loop introduces a selection effect, since a reported verdict may be the product of up to five attempts; in this corpus, however, twelve of thirteen runs closed without diagnosis-layer repair, the single exception being disclosed with the configuration (Section~\ref{sec:setup}). Third, the protocol's numeric thresholds---the $\ge$2-elimination floor, the 0.90 exclusion-confidence requirement, the 0.70/0.65/0.50 caps, the judge's 90-point gate---are fixed policy constants applied unchanged across the corpus; no threshold-sensitivity arm has been run, and the headline counts are conditional on them. The model-controlled ablation (Section~\ref{sec:ablation}) already bounds the accuracy claim (9/9 bare versus 4/9 resolved)---so the pipeline's value on documented faults is process quality, not accuracy. Expanding the repeats (all 12 scenarios, $k\ge3$: verdict-type stability, confidence spread, repair counts) remains pre-registered, with eleven scenarios awaiting their second in-era execution---gaps we report rather than fill; extending the corpus (the full 20-class TEP catalogue with the extended simulation dataset~\cite{reinartz2021extended}; bearing datasets; novel-disturbance suites) is the designed path to tighter intervals.

\textbf{Prompt-knowledge contamination.} TEP's documented fault mechanisms appear in textbooks and are therefore plausibly memorized by frontier models. Our protocol mitigates this at the level where memorization actually operates---but the mitigation has a limit. The documented cause table records not only the disturbed channel but each fault's dynamic form (IDV4 step, IDV11 random variation, IDV14 stiction), so recall of that table alone could reproduce the observed mechanism labels; the statistics-driven reading rests on the recorded reasoning chains, each of which cites the per-recording statistic that discriminates its form. Two residual channels are disclosed. First, the scenario identifier: every TEP \emph{and SKAB} case id and run-directory name embeds the documented mechanism (and dynamic form where the catalogue records one)---\texttt{tep\_d14\_reactor\_valve\_sticking} names the answer outright, and \texttt{skab\_cavitation\_13}, the corpus's only brief-blind scenario, is no exception---so identifier recall is possible in principle across the documented catalogue; IDV14, one of the four resolved verdicts, is the sharpest case. The leakage sentinel machine-checks brief and user-context content, not directory names, so the identifier channel sits outside it. A second, adjacent channel is the dispatcher itself: the twelve canonical scenarios were batched under one orchestrating session, whose context could in principle carry directory names across scenarios---the gates bound what stages write, not what the dispatcher has read---so the committed neutralized-identifier re-run mandates one orchestrating session per scenario as well.

The joint worst case is stated rather than implied: no resolved verdict survives both the identifier-recall discount on the TEP-derived outcomes and the identity-strict batch-93 discount---the four resolved verdicts are precisely three TEP hits plus batch 93---so the joint floor is 0/9, and the identity-strict reading alone is 3/9; in Table~\ref{tab:overall} and Fig.~\ref{fig:benchmark} the three TEP-derived resolved verdicts should be read as a recall-possible tier until the neutralized-identifier re-run executes. That re-run (the six TEP and two SKAB fault scenarios under opaque hash run identifiers and identifier-stripped brief headers, $\approx$16\,h of agent time by the cost accounting below) is committed as the first post-revision execution; it requires a contained harness extension (hash-id run directories) tracked in the released repository. The briefs themselves carry no fault vocabulary (leakage-sentinel checked), and the form-discrimination evidence of Section~\ref{sec:permech} is the main evidence-backed mitigation among the disclosed residual channels: the mechanism labels the pipeline commits to are read off per-recording statistics, which identifier recall would not require. Under the identifier channel alone, recall could supply both the mechanism and its dynamic form---sharpening, not replacing, the recall-possible reading of the three TEP-derived resolved verdicts (IDV1, IDV7, IDV14) stated above. As a partial counter-check, the SKAB ``other''-group scenario, whose fault label is withheld entirely from all natural-language brief content, still yielded the cavitation-type mechanism ranked among the survivors of a capped verdict---and the ranking offers weak counter-evidence against recall doing the work: the cap-internal order tracks five-factor evidence weights (Section~\ref{sec:calibration}), so a directory-supplied label with no recording evidence would not be expected to out-rank a statistically favored survivor---and cavitation ranks second of two, behind the valve-side hypothesis whose statistics the recording favors. Briefs are clean; directory names are not; both facts are disclosed. The RAG channel was unavailable in every run (Section~\ref{sec:arch-layers}), and no L6-tagged item exists anywhere in the released corpus: no verdict rests on external evidence---every result in this paper was obtained in the degraded, no-external-knowledge mode. A novel-disturbance suite (combination perturbations absent from the literature) is the definitive control and remains under construction.

\textbf{What the five capped verdicts mean.} The five \mbox{\textsc{competing\_set}} outcomes are the protocol working as designed, and they are also the benchmark's most transferable finding: each verdict localizes the disturbance to a subsystem, eliminates whole mechanism families with quantitative contradictions, keeps the true mechanism ranked among the survivors, and names the one measurement whose addition would collapse the set to \mbox{\textsc{determined}} (Section~\ref{sec:calibration}). In deployment terms these are not failures but work orders for targeted instrumentation---with one attribution caveat: a deterministic alternative path to the same content (localize the subsystem from the statistics, then enumerate that subsystem's unmeasured channels from the ontology) has not been tested by a script-only arm, so attributing this finding to the reasoning protocol rather than to sensor-set identifiability is a hypothesis; the UA-fork quantification of Section~\ref{sec:casestudy} is the current evidence that the reasoning layer adds more than a channel inventory. The three-state protocol is nonetheless what converts a diagnosis into that engineering workflow: \mbox{\textsc{determined}} conclusions carry mechanism chains that point to corrective action; capped conclusions carry the named channel. To keep the framing falsifiable, the released scenario definitions record the expected verdict set of every scenario, fixed before execution; SKAB inlet-valve throttling is the one scenario whose expectation was \mbox{\textsc{determined}}-only, and its capped verdict is counted as the benchmark's single compliance miss rather than narrated as a success.

\textbf{Operating semantics of capped verdicts.} A capped verdict is actionable only if its consumer knows what to do between the verdict and the named measurement's installation---an instrumentation lead time of weeks or months during which the anomaly is live. The named channel says what would settle the set, not what to do meanwhile, so a deployment needs three operating complements that the benchmark does not measure: (i)~an \emph{interim envelope} covering both survivors (operate to the intersection of the conditions under which each survivor is safe); (ii)~an \emph{escalation trigger} (the monitored signature whose drift past a stated bound re-opens or tightens the set); and (iii)~a \emph{decision-rights rule} mapping states to actions (\mbox{\textsc{determined}} $\to$ maintenance-approval path; \mbox{\textsc{competing\_set}} $\to$ deferred observation with explicit acceptance criteria; \textsc{needs\_data} $\to$ a data-acquisition task). The corpus's 5/9 capped base rate also matters for trust calibration: an operator half of whose verdicts are capped must sustain appropriate reliance (Section~\ref{sec:trust}) across both states, and under-capping---the costly direction the benchmark's one-directional typing cannot see (Section~\ref{sec:calibration})---is where that reliance would erode first. For IDV4 the envelope is concrete: operate the reactor temperature inside its pre-onset band with the cooling valve XMV\_10 below its pre-onset ceiling---conditions safe under both survivors---and treat sustained drift of the reactor-temperature residual beyond its pre-onset $\sigma$ as the escalation trigger that re-opens the set.

\textbf{Inference authorship and provenance.} Every scored artifact in this benchmark is authored by a gated pipeline stage---the diagnosis by the diagnostician, the quality verdict by the judge, the ENDORSED decision by the physical auditor, the review verdict by the independent HTML reviewer---and each carries a provenance label in the grading record. The runs reported here executed in the production dispatch form (independent per-stage sub-agent sessions; Section~\ref{sec:framework}), so the reviewing stages were physically separate sessions from the authoring stages---separate in session and protocol, though sharing the one model family (Section~\ref{sec:framework}). Because self-declared quality scores are not verifiable, the benchmark's reported quality metric is the deterministic compliance rubric (Section~\ref{sec:setup})---and the rubric, the grading keywords, and the allowed verdict sets remain author-designed instruments measuring conformance to this paper's own protocol; their release with the repository makes them auditable, but independent anchoring (the blinded expert adjudication of the released evidence chains) is the outstanding validation step---so the ``process quality'' margins claimed for the pipeline are architectural properties, not yet measured outcomes. A further provenance caveat: the L1--L7 levels are asserted by the reasoning stage itself; the machine gate checks ceiling arithmetic and grounding, not the correctness of a label, so a mislabeled evidence item could inflate its own ceiling---a disclosed residual risk whose mechanical binding (L1/L3 $\to$ ingestion/statistics files, L5 $\to$ ontology nodes) is a toolchain next step. An internal integrity audit that found (and removed) a scripted shortcut capable of fabricating artifacts and agent events is itself part of the released history; the completeness gate now makes that shortcut unreplayable. The audit plane's trust boundary is internal: the event log proves process conformance, not external attestation---a compromised execution layer could still fabricate events---which is why the grader's division of labour (scripts verify, never author) and the artifact cross-checks carry the weight. What is still missing is a comparison against \emph{human} experts on identical digests; an expert-vs-pipeline blind test is planned; and a mutation-based sensitivity audit of the gates themselves---injecting corrupted artifacts and wrong diagnoses to measure per-gate capture rates---is pre-registered alongside the repeat study. Table~\ref{tab:trustboundary} consolidates this trust model, stated piecemeal across Sections~\ref{sec:arch}, \ref{sec:framework}, and the present section: for each assumption the audit plane relies on, the mechanism that enforces it, the residual risk that remains, and the committed remedy.

\begin{table}[!t]
\caption{The audit plane's trust model, consolidated: each assumption, its enforcing mechanism, the residual risk disclosed in the text, and the committed remedy or status.}
\label{tab:trustboundary}
\centering
\footnotesize
\resizebox{\textwidth}{!}{%
\begin{tabular}{@{}p{2.6cm}p{4.4cm}p{4.6cm}p{4.0cm}@{}}
\toprule
Assumption & Enforcing mechanism & Residual risk (disclosed) & Remedy / status \\
\midrule
Truth isolation & bidirectional leakage sentinel over briefs and user-context files & run-directory names carry mechanism words (identifier channel, SKAB included) & neutralized-identifier re-run committed (eight fault scenarios) \\
Execution proof & event log + finalization gate + 14-artifact completeness set & pipeline-internal records, not externally signed attestations & hash-chained, append-only event log (toolchain next step) \\
Evidence levels & machine gate checks ceiling arithmetic and grounding & L1--L7 labels are asserted by the reasoning stage itself & mechanical binding L1/L3$\to$files, L5$\to$ontology (next step) \\
Reviewer independence & per-stage sub-agent sessions; judge is not a benchmark metric & one model family: correlated generator--reviewer failures possible & engine-substitution arm (pre-registered) \\
Scoring instrument & deterministic rubric, released with edit history & author-designed; limited discriminative power (9/12 at ceiling) & blinded expert adjudication (planned) \\
\bottomrule
\end{tabular}%
}
\end{table}

\textbf{Cost.} The clean per-scenario figure for the full gated pipeline is the sequential drawn re-execution: $\approx$2.0\,h of agent time end to end (117\,min event span), of which the deterministic statistical stage---script-speed, seconds---is a negligible share. Per-stage wall-clock spans across the thirteen released event logs give a median full-pipeline agent span of $\approx$1.9\,h ($\approx$112\,min: ontology $\approx$14\,min, statistics $\approx$22\,min, diagnosis $\approx$20\,min, the two review stages $\approx$29\,min combined, reporting and HTML the remaining $\approx$26\,min). The twelve canonical scenarios were executed as concurrent batches under one orchestrating agent (per-run event spans of $\approx$6\,h therefore reflect batching, not isolated effort; Section~\ref{sec:setup}). The pipeline is orders of magnitude more expensive than a fitted classifier at inference time and far slower than single-prompt baselines. The cost profile bounds the task fit: episodic, high-consequence investigations---batch-deviation, GxP, unit-trip RCA, where expert studies already run hours to days---are the natural regime; minute-scale alarm triage is out of scope by construction, and \textsc{needs\_data} is the designed answer when the recording cannot support the question. The economic comparison is only meaningful against the \emph{engineering cost of obtaining root causes} (expert hours), and the ontology-asset layer (Section~\ref{sec:arch-layers}) amortizes the reasoning cost across a plant's scenario fleet.

\subsection{Industrial Deployment Prospects}
We separate what the benchmark and one operational deployment \emph{demonstrate} from what the architecture \emph{projects}.

\textbf{One deployment anecdote.} This example is self-reported and unaudited (artifacts anonymised, not part of the released corpus): a film-line run exercised the pipeline end to end on episodic scratch defects of a biaxially-oriented polyester film line and returned \mbox{\textsc{needs\_data}} with an explicit discriminating-data list after 180 parameter--defect associations failed false-discovery-rate correction---no cause was fabricated. Projected: the integration, sovereignty, and documentation claims below are design properties of the architecture, not yet benchmarked outcomes.

\textbf{Integration with existing plant infrastructure.} IDD consumes what plants already produce: flat exports from historians (OPC-UA/PI), LIMS quality records, and batch logs. The same pipeline runs interactively through a web application and programmatically through a versioned REST API (\texttt{/api/diagnosis} with engine selection and typed error contracts), so a deployment is an export job plus API integration rather than an instrumentation project---the integration work is real but bounded by the historian export format. One export hazard deserves its own warning: historian compression (swinging-door and deadband algorithms) systematically destroys or fabricates exactly the statistics this pipeline's signatures rest on---variance ratios, quantised stepping, frozen means---so an export must request uncompressed or fixed-period data, and the export's compression settings belong in the run configuration and its fingerprint. The first site check for a missing channel is whether the tag exists but is unmapped in the export, which is more common than true absence. Because the deterministic stage is byte-reproducible, results do not drift when the surrounding plant IT systems change.

\textbf{Data-sovereign deployment.} The fourteen-engine execution layer (Section~\ref{sec:arch}) accepts on-premise and private-model agents through uniform dialect adapters. For operators who cannot send process data to external LLM services---defence, pharmaceuticals, semiconductors, national infrastructure---the reasoning stage can be pointed at an internally hosted model without touching the audit, evidence, or gate contracts---no non-default engine has yet executed the full gated pipeline in the released corpus, so this path is design-level, like the amortization projection; the benchmark protocol itself executes identically against a scripted mock engine, which is how the full chain is exercised in CI.

\textbf{Documentation for regulated operations.} Because every conclusion is bound to graded evidence, physics checks, and a fingerprinted execution trace, the by-product of each run is an RCA report whose traceability serves the documentation expectations of regulated environments (e.g., GxP batch-deviation investigation, in the sense of ALCOA+-style data integrity~\cite{mhra2018} and GAMP~5 computerised-systems thinking~\cite{gamp2022}). This is a documentation property of the artifact, not a computerised-system validation claim; the latter would require a formal qualification exercise.

\textbf{Deployment governance and knowledge compounding.} IDD is positioned as decision support: a \mbox{\textsc{determined}} verdict does not authorise maintenance action---a human engineer signs off. In the current release the sign-off is procedural (no separate human-endorsement artifact exists in the run contract, and all thirteen benchmark runs were fully autonomous); we disclose this as an implementation gap, not a deployed control, and the three-state protocol is precisely the escalation interface between machine conclusion and human decision. Model-change management rides the engine layer's version disclosure (Section~\ref{sec:setup}): a provider silently updating the reasoning model changes the system, so the benchmark configuration disclosure doubles as the deployment pinning record. A safety-context boundary is equally explicit: IDD is an off-line advisory layer for post-event investigation; it must not be connected to any protection or interlock path---in the sense of safety instrumented functions (IEC~61511, \emph{Functional safety---Safety instrumented systems for the process industry sector})---and the three-state escalation interface presupposes a human decision-maker outside the loop. The system's own model card (model family and snapshot, protocol version, known failure modes: identifier recall, one-directional verdict typing, the two cap sets' gap) ships with the repository. And because the evaluation harness and this manuscript were produced with AI assistance inside the same human--AI loop that produced the system, independent deployment validation is a precondition for safety-relevant use. Finally, the ontology layer compounds knowledge across runs (Section~\ref{sec:arch-layers}): all thirteen benchmark runs published their ontologies to the store on completion, and every field incident becomes a candidate graded scenario, closing the loop between deployment and evaluation. In deployment the marginal cost of a diagnosis therefore falls as the fleet grows (projected: the benchmark pinned full rebuilds, so this amortization is design-level, not measured)---an economics that sample-level detectors, which must be retrained per asset, do not share.

\textbf{Scope of applicability.} The method requires multivariate time-ordered process data with engineering semantics for the channels (to build the ontology) and applies to continuous, batch, and testbed processes alike---the three benchmark families were chosen to span exactly that range. It does not apply where no time-aligned process data exists, and its root-cause reach is bounded by the sensor set, as the capped verdicts demonstrate.

\section{Conclusion}
\label{sec:conclusion}

We presented IDD, an agentic diagnosis pipeline that turns LLM-based root-cause analysis into an auditable engineering procedure---competing hypotheses with mechanism classes, evidence levels bounding reported confidence, a four-condition anti-spurious filter, physical verification, and quality gates with bounded repair---together with a machine-checked evaluation contract for diagnosis benchmarks: truth isolation, data fingerprinting, stage-authored artifacts with provenance, and gate-graded scoring. On the 12-scenario corpus the headline outcome is deliberately modest: 4/9 faults resolved Top-1, the true mechanism ranked in every fault scenario's surviving set (top-$k$ 9/9; 8/9 when a survivor must carry the specific form or stream branch, not just the family)---a family-level set-membership statement, not a specificity-tested resolution bound---zero false alarms (one of three blind), five capped verdicts that each name the measurement that would settle them, and every number backed by an execution proof. The model-controlled ablation bounds what the pipeline can claim: a bare call to the same model already exceeds its accuracy (9/9 versus 4/9), so the contribution is the auditable process rather than the accuracy. One seeded re-execution reproduced verdict state, actuator, and mechanism class within two confidence points; the deterministic suite reproduced byte-for-byte; and the per-fault TEP comparison shows IDD resolving the GPT-4o baseline's IDV1 miss while scoring the PCA-excluded pair that baselines exclude, as capped verdicts with the mechanism family ranked.
These claims have limits: resolution of documented faults is the pipeline's weak axis by its own ablation (a bare same-model call scores 9/9 versus the pipeline's 4/9 resolved; accuracy is not the contribution), results are single recorded executions of a stochastic reasoning layer whose configuration is disclosed and whose run identifiers carry the documented mechanism labels (so no headline verdict is recall-clean until the committed neutralized-identifier re-run executes), the machine gate's flat 0.90 non-CS ceiling flags two verdicts (0.91 and 0.93) that the agent-side cap set permits, so the two cap sets' gap is measured rather than hidden; and the rubric and keywords are author-designed instruments with independent expert anchoring outstanding. The evidence base is one model family at one deployment on twelve public-dataset scenarios; cross-model and cross-plant generalization of the process guarantees (in contrast to the accuracy findings, which the ablation already bounds) remains to be demonstrated.
The evaluation framework, per-scenario evidence chains, and consistency artifacts are released to support scrutiny and reuse; the pre-registered repeat and ablation studies, and the novel-disturbance suite, are the immediate next steps.

\section*{CRediT authorship contribution statement}
\textbf{Haozheng Zhang}: Conceptualization, methodology, software, investigation, data curation, writing -- original draft, writing -- review \& editing.

\section*{Declaration of competing interest}
The authors declare that they have no known competing financial interests or personal relationships that could have appeared to influence the work reported in this paper.

\section*{Funding}
This research did not receive any specific grant from funding agencies in the public, commercial, or not-for-profit sectors.

\section*{Data availability}
The system and harness source is available at \url{https://github.com/kingdol666/industrial-deep-diagnostic}; the complete evaluation artifact layer accompanies this submission as supplementary material, and an archived snapshot with a persistent identifier (Zenodo DOI) will be minted on acceptance. The artifact layer comprises: the frozen scenario definitions with grading keyword sets (pinned at commit \texttt{221d061}) and expected verdict sets; per-fault literature baselines with edit history; per-scenario blind inputs and briefs with SHA-256 fingerprints (\texttt{baselines\_input\_audit.json}); the thirteen pipeline run directories, whose stage-authored artifacts include the reasoning chains (\texttt{04\_diagnostics}), the judge feedback (\texttt{05\_review} judge records), and append-only execution-event logs that carry the per-scenario repair accounting; the rubric (\texttt{rubric.json}) and grading records with provenance labels; and the aggregation, consistency-verification, and report scripts with the Markdown/HTML scoring reports. The deposit contains exactly the thirteen canonical run directories (the recorded 2026-09-17 batch and the drawn re-execution); an interrupted, never-graded second-round batch on the workspace is excluded from the deposit and from every statistic in this paper. Third-party datasets are publicly available from their cited sources.

\section*{Declaration of generative AI and AI-assisted technologies in the writing process}
During the preparation of this work the author used AI-assisted coding and writing tools (the GLM-based agent CLI used throughout the project's harness, and editor-level writing assistants) in order to develop the evaluation harness, execute the benchmark pipeline, and support manuscript drafting and editing. After using these tools, the author reviewed and edited the content as needed and takes full responsibility for the content of the published article. The agentic system that is the \emph{subject} of this paper (IDD) is unrelated to the writing-assistance tools used during preparation; all benchmark executions reported here were produced by the IDD pipeline itself under the released evaluation harness.

\bibliographystyle{elsarticle-num}
\bibliography{refs}

\appendix

\section{Supplementary Tables and Released-Artifact Index}\label{app:artifacts}

\setcounter{table}{0}
\renewcommand{\thetable}{A\arabic{table}}
\renewcommand{\theHtable}{A\arabic{table}}

Table~\ref{tab:perfault} transcribes the per-fault TEP comparison discussed in
Section~\ref{sec:comparison}. The artifacts released for every benchmark scenario are exactly
the fourteen-item completeness set of Table~\ref{tab:artifacts}, laid out one run directory
per scenario in the released repository, each machine-readable and carrying the execution
proofs described in Section~\ref{sec:framework}. The machine-generated English benchmark
report is \texttt{results/benchmark/benchmark\_report\_en.md}; its own appendix lists the
grader-side ground truth and must not be exposed to a diagnosing agent.

\begin{table}[!t]
\caption{Per-Fault Outcomes on the TEP Subset (machine-readable companion to Fig.~\ref{fig:perfault}). FE = FaultExplainer (root-causes-included prompt, top-3 candidate list, alias hits accepted); ``unscored'' = excluded from FE scoring by its PCA implementation~\cite{chiang2001fault}. IDD: strict mechanism-keyword scoring, no candidate list; DET = \mbox{\textsc{determined}}; CS = capped \mbox{\textsc{competing\_set}} (every CS verdict ranks the true mechanism among its survivors; TEP-subset top-$k$ = 6/6, or 5/6 under the branch-localized reading of Section~\ref{sec:setup}, which counts IDV3's stripper-branch membership as a miss). In FE's additional no-candidate regime (General-Reasoning prompt, top-3, lenient ``related'' grading), both models score 8/11---the reported figure, which FE's own text contradicts (it names faults 10 and 13 as failing); we transcribe it as reported, not as an independent count. IDD's Top-1 summary: 3/6 (Wilson CI $[18.8\%,81.2\%]$, $n{=}6$); shared detectable subset 3/4.}
\label{tab:perfault}
\centering
\resizebox{\textwidth}{!}{%
\begin{tabular}{llccp{10.5cm}}
\toprule
TEP fault & FE GPT-4o & FE o1-preview & PCA detectable & IDD verdict (conf.) \\
\midrule
IDV1 A/C feed-ratio step (Stream 4) & wrong & correct & yes & \textbf{Top-1 correct} DET 0.91 \\
IDV3 D-feed temperature step & unscored & unscored & \textbf{no} & CS 0.55 (steam-loop site; fork named --- survivor names ``D/E feed temperature or stream-10 enthalpy'', the truth family inside a stripper-feed-branch hypothesis) \\
IDV4 reactor cooling-water temp.\ (step) & unscored & unscored & \textbf{no} & CS 0.63 (heat-removal step top-ranked; inlet/UA fork named) \\
IDV7 C-header pressure (Stream 4) & correct & correct & yes & \textbf{Top-1 correct} DET 0.78 \\
IDV11 cooling-water temp.\ (random) & correct & correct & yes & CS 0.60 (random-variation form; inlet/UA fork named) \\
IDV14 cooling-water valve (stiction) & correct & correct & yes & \textbf{Top-1 correct} DET 0.80 \\
\midrule
Scoreable-set summary & 7/11 & 9/11 &---& 3/6 (CI $[18.8\%,81.2\%]$, $n{=}6$); shared detectable subset 3/4; top-$k$ 6/6 \\
\midrule
FE no-candidate regime (top-3, lenient) & \multicolumn{2}{c}{8/11 both models} &---&---\\
\bottomrule
\end{tabular}%
}
\end{table}

Table~\ref{tab:baselinedetail} completes the per-scenario record for the two model-controlled baseline arms of Section~\ref{sec:ablation}: the bare single-call LLM's actual top-1 answer and the classical PCA arm's SPE-contribution top-3 variables, alongside the IDD verdict.

\begin{table}[!t]
\caption{Per-Scenario Baseline-Arm Content (companion to Table~\ref{tab:ablation} and Fig.~\ref{fig:matrix}; every entry read from the released records \texttt{baselines.json} and \texttt{baseline\_pca\_rca.json}). Bare-LLM top-1 answers are abbreviated English glosses of the released verbatim strings (one is in Chinese); ``(hit)'' = strict mechanism-keyword hit. PCA columns are the SPE-contribution top-3 variables---variable-level rankings, not mechanism verdicts. $\dagger$ = control scenario at the nominal $\approx$1\% false-alarm level (contribution lists on nominal rows are not diagnostic). The bare arm covered two of three controls (TEP control not run bare). IndPenSim abbreviations: Fh = heating-water flow, Q = generated heat, Fc = cooling-water flow, Fs = sugar-feed flow, Fb = base flow. TEP channels: XMEAS\_2/9/21 = D-feed flow, reactor temperature, reactor cooling-water outlet temperature; XMEAS\_25/29/31/33/38 = Downs--Vogel stream-composition analyzer channels; XMV\_n = $n$-th manipulated variable. IDD verdicts repeated from Table~\ref{tab:mech} for side-by-side reading. SHA-256 fingerprints of the baseline briefs are released in \texttt{baselines\_input\_audit.json}.}
\label{tab:baselinedetail}
\centering
\scriptsize\setlength{\tabcolsep}{3.4pt}\renewcommand{\arraystretch}{0.94}
\resizebox{\textwidth}{!}{%
\begin{tabular}{@{}p{3.0cm}p{5.6cm}p{4.6cm}p{2.6cm}@{}}
\toprule
Scenario & Bare-LLM top-1 answer (strict) & PCA SPE top-3 contributions & IDD verdict (conf.) \\
\midrule
SKAB inlet-valve throttling & ``Valve malfunction---partial closure/sticking of the loop valve causing flow restriction'' (hit) & Accel2RMS $\cdot$ Accel1RMS $\cdot$ FlowRMS & CS 0.60 \\
SKAB cavitation (blind) & ``Centrifugal pump cavitation: suction-side pressure deficit $\to$ bubble formation and collapse'' (hit) & FlowRMS $\cdot$ Accel2RMS $\cdot$ Accel1RMS & CS 0.60 \\
Control: SKAB anomaly-free & normal (no fault)$^\dagger$ & FlowRMS $\cdot$ Accel2RMS $\cdot$ Accel1RMS$^\dagger$ & normal 0.86 \\
\midrule
TEP IDV1 A/C feed-ratio step & ``IDV(1): A/C feed ratio step change in Stream 4 (B composition constant)'' (hit) & XMEAS\_21 $\cdot$ XMEAS\_2 $\cdot$ XMEAS\_38 & DET 0.91 \\
TEP IDV3 D-feed temp.\ (hard) & ``IDV(3): step change in D feed temperature (stream 2)'' (hit) & XMEAS\_31 $\cdot$ XMEAS\_38 $\cdot$ XMEAS\_29 & CS 0.55 \\
Control: TEP fault-free & --- (arm not run) & XMEAS\_38 $\cdot$ XMEAS\_29 $\cdot$ XMEAS\_25$^\dagger$ & normal 0.90 \\
TEP IDV4 cooling (step) & ``IDV(4): reactor cooling-water inlet temperature step change'' (hit) & XMV\_10 $\cdot$ XMEAS\_9 $\cdot$ XMEAS\_31 & CS 0.63 \\
TEP IDV7 C-header pressure & ``IDV(7): C header pressure reduced (step change)'' (hit) & XMEAS\_33 $\cdot$ XMEAS\_29 $\cdot$ XMEAS\_38 & DET 0.78 \\
TEP IDV11 cooling (random) & ``IDV(11): random variation in reactor cooling-water inlet temperature'' (hit) & XMV\_10 $\cdot$ XMEAS\_9 $\cdot$ XMEAS\_31 & CS 0.60 \\
TEP IDV14 valve stiction & ``IDV(14): reactor cooling-water valve sticking'' (hit) & XMEAS\_21 $\cdot$ XMEAS\_2 $\cdot$ XMV\_1 & DET 0.80 \\
\midrule
IndPenSim batch-93 deviation & ``pH control-loop disturbance---acid-feed ($F_a$) extreme excursions, consistent with acid-pump malfunction or pH-probe drift'' (hit) & Fh $\cdot$ Q $\cdot$ Fc & DET 0.82 \\
Control: IndPenSim batch-1 & normal (no fault)$^\dagger$ & Fc $\cdot$ Fs $\cdot$ Fb$^\dagger$ & normal 0.93 \\
\bottomrule
\end{tabular}%
}
\end{table}

\end{document}
